\section{Proofs for point output}\label{app:points}

This appendix proves the results of Section~\ref{sec:points}.
Section~\ref{app:points-tools} collects the tools: explicit sparse
evaluation, interpolation inequalities, integer eigenvalue and Hoffman
bounds, relative coordinates, the convex value interface, and the selector
lemma. Sections~\ref{app:points-global} and~\ref{app:points-cubic} prove
Theorems~\ref{thm:global-point} and~\ref{thm:cubic-point}.
Section~\ref{app:points-boundary} contains the constructions behind the
boundary results of Sections~\ref{sec:points-limits}
and~\ref{sec:points-representation}.

Throughout, $\langle r\rangle$ denotes the binary encoding length of a
rational number, vector, or matrix, written with integer numerators and
positive integer denominators. The norm $\norm{\cdot}$ is Euclidean, and
for a matrix it is the spectral norm. For a matrix $T$ we write
$\norm{T}_\triangle=\sum_{i,j}|T_{ij}|$; it bounds the spectral norm of $T$
and of $T^{\mathsf T}$. Running times count bit operations. A bound
``polynomial in $a,b,\ldots$'' is a polynomial in $a+b+\cdots$ with an
absolute exponent.

\subsection{Tools}\label{app:points-tools}

\paragraph{Classical rational subroutines.}
Three classical facts are used without proof. First, Gaussian elimination
over $\Q$ computes ranks, bases of kernels and ranges, inverses, and
solutions of linear systems in polynomial time, and all its outputs have
polynomial encoding length \citep[Section~1.4]{GroetschelLovaszSchrijver1988}.
Second, rational linear programs are solvable in polynomial time: given
rational $A,b,c$, one can decide whether $\{x:Ax\le b\}$ is empty and
whether $c^{\mathsf T}x$ is unbounded above on it, and otherwise compute an
optimal solution of polynomial encoding length; systems that also contain
linear equations are handled in the same way
\citep[Theorem~(6.4.12) and Section~6.5]{GroetschelLovaszSchrijver1988}.
Third, Farkas' lemma: for real matrices $G_1,G_2$ and a real vector $w$,
the system $G_1y_1+G_2y_2=w$, $y_1\ge0$, has no solution if and only if
some $d$ satisfies $G_1^{\mathsf T}d\le0$, $G_2^{\mathsf T}d=0$, and
$w^{\mathsf T}d>0$. With rational data such a $d$ with $w^{\mathsf T}d=1$
is an optimal solution of the linear program
$\max\{w^{\mathsf T}d:G_1^{\mathsf T}d\le0,\ G_2^{\mathsf T}d=0,\
w^{\mathsf T}d\le1\}$.

\paragraph{Explicit sparse evaluation.}
An explicit sparse polynomial is a list of pairs $(f_\alpha,\alpha)$ with
distinct exponent vectors $\alpha$ and nonzero rational coefficients. Its
numerical degree is the largest $|\alpha|=\sum_i\alpha_i$ in the list.
Exponents may be written in binary; the bounds below are polynomial in the
numerical degree, which is polynomial in the input length exactly when the
degree is.

\begin{lemma}[Explicit sparse evaluation]\label{lem:points-sparse}
Let $f=\sum_{\alpha\in\mathcal A}f_\alpha y^\alpha\in\Q[y_1,\ldots,y_N]$
be an explicit sparse polynomial of encoding length $L_f$ and numerical
degree at most $D\ge1$. Let $T\in\Q^{N\times n}$, $t\in\Q^N$, and
$\tau\in\Q$ with $\tau\ge0$, and put
\[
 \varphi(x)=f(Tx+t)+\tau\norm x^2 .
\]
\begin{enumerate}[label=(\alph*)]
\item For every rational $x$, the exact rational values $\varphi(x)$ and
$\nabla\varphi(x)$ can be computed in time polynomial in
$L_f,D,\langle T,t,\tau\rangle,\langle x\rangle$. The computation does not
form the expanded polynomial $f(Tx+t)$.
\item Let $\varrho\ge1$ be rational, put
$\varrho_1=\max\{1,\ \varrho\max_i\sum_j|T_{ij}|+\max_i|t_i|\}$, and
\[
 V_f=\sum_{\alpha}|f_\alpha|\varrho_1^{|\alpha|},\qquad
 G_f=\sum_{|\alpha|\ge1}|f_\alpha|\,|\alpha|\,\varrho_1^{|\alpha|-1},\qquad
 H_f=\sum_{|\alpha|\ge2}|f_\alpha|\,|\alpha|^2\varrho_1^{|\alpha|-2}.
\]
Then on the box $[-\varrho,\varrho]^n$,
\[
 |\varphi|\le V_f+\tau n\varrho^2,\qquad
 \norm{\nabla\varphi}\le\norm T_\triangle G_f+2\tau n\varrho,\qquad
 \norm{\nabla^2\varphi}\le\norm T_\triangle^2H_f+2\tau .
\]
These bounds are rationals of encoding length polynomial in
$L_f,D,\langle T,t,\tau,\varrho\rangle$.
\end{enumerate}
\end{lemma}

\begin{proof}
(a) Compute $y=Tx+t$. If every coordinate of $y$ has encoding length at most
$\sigma$, then each power $y_i^{\alpha_i}$, and hence each monomial
$y^\alpha$, is a rational whose numerator and denominator have at most
$D\sigma$ bits, computed by repeated multiplication. Summing the at most
$L_f$ terms $f_\alpha y^\alpha$ gives $f(y)$ with polynomially many bits.
The partial derivatives satisfy
$\partial_if(y)=\sum_{\alpha_i>0}\alpha_if_\alpha y^{\alpha-e_i}$ and are
computed in the same way. Finally
$\nabla\varphi(x)=T^{\mathsf T}\nabla f(y)+2\tau x$.

(b) For $\norm x_\infty\le\varrho$ we have
$\norm{Tx+t}_\infty\le\varrho_1$. For $\norm y_\infty\le\varrho_1$ and
$\varrho_1\ge1$, a monomial satisfies $|y^\alpha|\le\varrho_1^{|\alpha|}$,
the entries of its gradient satisfy
$\sum_i|\alpha_iy^{\alpha-e_i}|\le|\alpha|\varrho_1^{|\alpha|-1}$, and
the entries of its Hessian satisfy
$\sum_{i,j}|\alpha_i(\alpha_j-\delta_{ij})y^{\alpha-e_i-e_j}|\le
|\alpha|^2\varrho_1^{|\alpha|-2}$. The Euclidean norm of a vector and the
spectral norm of a matrix are bounded by the sums of the absolute values of
their entries. The chain rule gives
$\nabla_x f(Tx+t)=T^{\mathsf T}\nabla f(y)$ and
$\nabla^2_xf(Tx+t)=T^{\mathsf T}\nabla^2f(y)T$. The regularizer contributes
at most $\tau n\varrho^2$, $2\tau\norm x\le2\tau n\varrho$, and $2\tau$.
Each displayed bound is a sum of at most $L_f+1$ rationals whose lengths
are polynomial in the stated quantities.
\end{proof}

\paragraph{Interpolation inequalities.}
For an integer $D\ge1$ put
\[
 c_0(D)=(D+1)D^D,\qquad c_1(D)=(D+1)D^{D+1},\qquad c_2(D)=(D+1)D^{D+2},
\]
and for an integer $k\ge1$ put $\beta_k=(k+1)2^kk^k$. These integers have
$O(D\log(D+1))$ and $O(k\log(k+1))$ bits.

\begin{lemma}[Interpolation]\label{lem:points-interp}
Let $r$ be a real univariate polynomial of degree at most $D\ge1$.
\begin{enumerate}[label=(\alph*)]
\item $|r(t)|\le c_0(D)(1+|t|)^D\max_{0\le j\le D}|r(j/D)|$ for every
real $t$.
\item If $\varrho>0$ and $|r(t)|\le W$ for $0\le t\le\varrho$, then
$|r'(0)|\le c_1(D)W/\varrho$.
\item $|r''(0)|\le c_2(D)\max_{0\le j\le D}|r(j/D)|$.
\item Let $k\ge1$ and let $p\in\R[z_1,\ldots,z_n]$ have degree at most $k$
in each variable. If $|p(z)|\le V$ for every $z\in[0,1]^n$, then every
coefficient of $p$ has absolute value at most $\beta_k^nV$.
\end{enumerate}
\end{lemma}

\begin{proof}
For $j=0,\ldots,D$ let
$\lambda_j(t)=\prod_{i\ne j}(t-i/D)/\prod_{i\ne j}(j/D-i/D)$ be the
Lagrange basis for the nodes $j/D$, so that
$r=\sum_jr(j/D)\lambda_j$. The denominator of $\lambda_j$ has absolute
value $j!(D-j)!/D^D\ge D^{-D}$.

(a) Every numerator factor satisfies $|t-i/D|\le1+|t|$, so
$|\lambda_j(t)|\le D^D(1+|t|)^D$. Sum over the $D+1$ nodes.

(b) The derivative at zero of $\prod_{i\ne j}(t-i/D)$ is a sum of $D$
products of $D-1$ numbers in $[-1,0]$, so it has absolute value at most
$D$, and $|\lambda_j'(0)|\le D^{D+1}$. Apply this to
$s\mapsto r(\varrho s)$, whose values at the nodes lie in $[-W,W]$ and whose
derivative at zero is $\varrho r'(0)$.

(c) The second derivative at zero of the numerator is a sum over the
$D(D-1)$ ordered pairs of distinct indices of products of numbers in
$[-1,0]$. Hence $|\lambda_j''(0)|\le D(D-1)D^D\le D^{D+2}$.

(d) Use the nodes $j/k$ in $[0,1]$ and the corresponding univariate Lagrange
polynomials $\mu_0,\ldots,\mu_k$. The sum of the absolute values of the
coefficients is submultiplicative, and for a factor $z-i/k$ it equals
$1+i/k\le2$; with the denominator bound, each $\mu_j$ has coefficient sum at
most $2^kk^k$. Tensor interpolation
\[
 p(z)=\sum_{J\in\{0,\ldots,k\}^n}p(J/k)\prod_{i=1}^n\mu_{J_i}(z_i)
\]
is exact for polynomials of degree at most $k$ in each variable. The
coefficient sum of each product is at most $(2^kk^k)^n$, and there are
$(k+1)^n$ terms. Hence the coefficient sum of $p$ is at most $\beta_k^nV$.
The interpolation grid enters only this proof; no algorithm below
evaluates $p$ on it.
\end{proof}

\paragraph{Integer matrices.}

\begin{lemma}[Positive eigenvalues of integer matrices]\label{lem:points-integer-eigen}
Let $N$ be a symmetric positive semidefinite integer matrix of order $n$,
rank $r\ge1$, and entries bounded in absolute value by $C\ge1$. Every
nonzero eigenvalue of $N$ is at least $(nC)^{-(r-1)}$. Consequently, if
$M\succeq0$ is rational, $Q_M$ is a positive integer with $Q_MM$ integral,
$C\ge1$ bounds the entries of $Q_MM$, and $\Pi$ is the orthogonal
projection onto the range of $M$, then
$x^{\mathsf T}Mx\ge\norm{\Pi x}^2/(Q_M(nC)^{n-1})$ for every $x$.
\end{lemma}

\begin{proof}
The product of the nonzero eigenvalues of $N$ equals, up to sign, a
coefficient of its characteristic polynomial; it is the sum of the
principal minors of order $r$, an integer. It is positive, so it is at
least one. Each eigenvalue is at most the largest absolute row sum, which
is at most $nC$. Hence the smallest nonzero eigenvalue is at least
$(nC)^{-(r-1)}$. For the consequence, $Q_MM$ and $M$ have the same range,
$x^{\mathsf T}Mx=(\Pi x)^{\mathsf T}M(\Pi x)$, and $\Pi x$ lies in the
span of the eigenvectors of $Q_MM$ with nonzero eigenvalues. If $M=0$,
both sides vanish.
\end{proof}

\begin{proof}[Proof of Lemma~\ref{lem:integer-hoffman}]
Let $x\in Q$ and let $y$ be the Euclidean projection of $x$ onto the
nonempty closed convex set $Z$. If $x=y$ there is nothing to prove. Let
$I$ be the set of inequality rows active at $y$. The vector $x-y$ lies in
the normal cone of the polyhedron $Z$ at $y$: the linear function
$(x-y)^{\mathsf T}z$ is maximized over $Z$ at $z=y$. By linear programming
duality this normal cone is
$\{A_I^{\mathsf T}\beta+T^{\mathsf T}\alpha:\beta\ge0\}$, so
$x-y=A_I^{\mathsf T}\beta+T^{\mathsf T}\alpha$ for some $\beta\ge0$ and some
$\alpha$.

We pass to a linearly independent representation. Choose rows of $T$ that
form a basis of its row space; the representation can be rewritten with
equality coefficients supported on these rows. If the chosen equality rows
together with the inequality rows of positive coefficient are linearly
dependent, some nontrivial relation among them gives a positive coefficient
to at least one inequality row, because the chosen equality rows are
independent. Subtracting the largest multiple of the relation that keeps
all inequality coefficients nonnegative removes at least one inequality
row. After finitely many steps
$x-y=R^{\mathsf T}\binom{\alpha'}{\beta'}$, where $R$ stacks linearly
independent rows $T'$ of $T$ and $A'$ of $A_I$, with $\beta'\ge0$. Let $k'$
be the number of rows of $R$.

By the Cauchy--Binet formula, $\det(RR^{\mathsf T})$ is the sum of the
squares of the $k'\times k'$ minors of $R$, so it is at least $\mu^2$. The
largest eigenvalue of $RR^{\mathsf T}$ is at most
$\norm R_1\norm R_\infty\le(k'C)(sC)\le(sC)^2$. Therefore the smallest
eigenvalue is at least $\mu^2(sC)^{-2(k'-1)}\ge\mu^2(sC)^{-2(s-1)}$, and
$\sigma_{\min}(R)\ge\mu(sC)^{-(s-1)}$. Since $R$ has full row rank,
\[
 \Bigl\|\binom{\alpha'}{\beta'}\Bigr\|\le\frac{\norm{x-y}}{\sigma_{\min}(R)}.
\]
Now
\[
 \norm{x-y}^2=\alpha'^{\mathsf T}T'(x-y)+\beta'^{\mathsf T}A'(x-y).
\]
Each row of $A'$ is active at $y$ and satisfied at $x\in Q$, so
$A'(x-y)\le0$ and the second term is nonpositive. Since $T'y=t'$, the first
term equals $\alpha'^{\mathsf T}(T'x-t')$, which is at most
$\norm{\alpha'}\,\norm{Tx-t}$. Combining,
$\norm{x-y}^2\le\norm{x-y}\,\norm{Tx-t}\,(sC)^{s-1}/\mu$. Divide by
$\norm{x-y}$.

For integer matrices, a linearly independent family of $k'$ rows has rank
$k'$, so some $k'\times k'$ minor is a nonzero integer and has absolute value
at least one.
\end{proof}

The proof uses neither boundedness of $Q$ nor rationality of $b,t$. When
the equality rows are absent, $Z=Q$ and the bound reads $\dist(x,Z)=0$.

\paragraph{Relative coordinates.}

\begin{lemma}[Relative coordinates and an inner ball]\label{lem:points-relative}
Let $P=\{x\in\R^n:Ax\le b\}$ be nonempty, with rational $A,b$, and let
$\varrho\ge1$ be rational with $P\subseteq[-\varrho,\varrho]^n$. In
polynomial time one can compute an integer $m$ with $0\le m\le n$, a point
$x_0\in P\cap\Q^n$, a matrix $V\in\Q^{n\times m}$ of rank $m$ that contains
the $m\times m$ identity matrix as a submatrix, and, when $m\ge1$, a
rational system $Gu\le h$ with nonzero rows $g_i^{\mathsf T}$ and $h>0$ such
that, with $P'=\{u\in\R^m:Gu\le h\}$,
\begin{enumerate}[label=(\roman*)]
\item $P=\{x_0+Vu:u\in P'\}$ (for $m=0$, $P=\{x_0\}$);
\item the ball of radius $r_0=\min_ih_i/\norm{g_i}_1$ about $0$ lies in
$P'$, and $\norm u\le R_0=2m\varrho$ for every $u\in P'$.
\end{enumerate}
All outputs have encoding length polynomial in $\langle A,b,\varrho\rangle$.
In particular $m=\dim P$.
\end{lemma}

\begin{proof}
For each row $a_i^{\mathsf T}x\le b_i$ solve
$\max\{b_i-a_i^{\mathsf T}x:x\in P\}$. Call the row an implicit equality
if the optimum is zero; otherwise record an optimal point $x^{(i)}$, which
has positive slack in row $i$. Let $x_0$ be the average of the recorded
points, or any rational point of $P$ if none was recorded. Then $x_0\in P$
and $x_0$ has positive slack in every row that is not an implicit equality.

Let $Ex=e$ be the implicit equalities. Gaussian elimination brings this
system to reduced row echelon form, with free coordinates
$F\subseteq\{1,\ldots,n\}$, $|F|=m$. Its solution set is
$\{x_0+Vu:u\in\R^m\}$, where $V$ has the identity matrix in the rows
indexed by $F$ and $u=(x-x_0)_F$. Every point of $P$ satisfies $Ex=e$, so
(i) holds with $P'=\{u:x_0+Vu\in P\}$.

Substitute $x=x_0+Vu$ into $Ax\le b$. An implicit equality row becomes
$0\le0$: it is constant on the solution set of $Ex=e$ and vanishes at
$x_0$. Any other row with $V^{\mathsf T}a_i=0$ is constant on that set and
has positive slack at $x_0$, so it is redundant. Delete these rows. The
remaining rows read $g_i^{\mathsf T}u\le h_i$ with $g_i=V^{\mathsf T}a_i\ne0$
and $h_i=b_i-a_i^{\mathsf T}x_0>0$. If $\norm u\le r_0$, then
$g_i^{\mathsf T}u\le\norm{g_i}\,\norm u\le\norm{g_i}_1r_0\le h_i$. For
$u\in P'$ we have $u=(x-x_0)_F$ with $x\in P$, so
$\norm u\le\norm u_1\le2m\varrho$. If $m\ge1$, boundedness of $P'$ forces at
least one remaining row. All steps are linear programs and Gaussian
elimination. The ball in $P'$ shows that $P'$ is full dimensional, hence
$\dim P=m$.
\end{proof}

\paragraph{The convex value interface.}
We use the following theorem of Gr\"otschel, Lov\'asz, and Schrijver. For
a convex body $\mathcal K\subseteq\R^d$ (a compact convex set with nonempty
interior) and $\varepsilon>0$, let $S(\mathcal K,\varepsilon)$ be the set
of points within distance $\varepsilon$ of $\mathcal K$, and let
$S(\mathcal K,-\varepsilon)$ be the set of points whose closed
$\varepsilon$-ball lies in $\mathcal K$. A \emph{weak separation oracle}
for $\mathcal K$ receives a rational $y$ and a rational $\delta>0$ and
either asserts $y\in S(\mathcal K,\delta)$ or returns a rational $c$ with
$\norm c_\infty=1$ and $c^{\mathsf T}z\le c^{\mathsf T}y+\delta$ for all
$z\in S(\mathcal K,-\delta)$.

\begin{lemma}[Weak optimization from weak separation;
\citealp{GroetschelLovaszSchrijver1988}]\label{lem:points-gls}
There is an oracle algorithm that, given $d$, a rational $R>0$ with
$\mathcal K\subseteq\{z:\norm z\le R\}$, a rational vector $c$, a rational
$\varepsilon>0$, and a weak separation oracle for $\mathcal K$, either
returns a rational $y\in S(\mathcal K,\varepsilon)$ with
$c^{\mathsf T}z\le c^{\mathsf T}y+\varepsilon$ for all
$z\in S(\mathcal K,-\varepsilon)$, or asserts that
$S(\mathcal K,-\varepsilon)$ is empty. If the oracle is realized by an
algorithm whose running time is polynomial in the encoding lengths of its
query and of a fixed description of $\mathcal K$, the total running time is
polynomial in $d$, $\langle R\rangle$, $\langle c\rangle$,
$\langle\varepsilon\rangle$, and the length of that description.
\end{lemma}

This is the weak optimization problem \citep[(2.1.10)]{GroetschelLovaszSchrijver1988}
for the circumscribed convex body $(\mathcal K;d,R)$ given by a weak
separation oracle \citep[(2.1.13), (2.1.16)]{GroetschelLovaszSchrijver1988},
solved in oracle-polynomial time by
\citet[Corollary~(4.2.7)]{GroetschelLovaszSchrijver1988}. The running-time
consequence is their composition rule for oracle-polynomial algorithms with
polynomially bounded oracle answers
\citep[(1.2.1) and Section~4.1]{GroetschelLovaszSchrijver1988}; we include
the description of $\mathcal K$ in every oracle query.

\begin{proof}[Proof of Lemma~\ref{lem:convex-value}]
(a) Decide by linear programming whether $P$ is empty. If it is not, apply
Lemma~\ref{lem:points-relative}. If $m=0$, return $\hat x=x_0$ and
$\ell=\varphi(x_0)$. Assume $m\ge1$ and let $\psi(u)=\varphi(x_0+Vu)$ on
$P'$; it is continuous and convex. For $u\in P'$ the point $x_0+Vu$ lies in
$P$, so $\psi(u)$ and the vector
$s_\psi(u)=V^{\mathsf T}s(x_0+Vu)$ are computable in time polynomial in
$\Delta+\langle u\rangle$. The vector $s_\psi(u)$ is a subgradient of $\psi$
on $P'$, that is, $\psi(u')\ge\psi(u)+s_\psi(u)^{\mathsf T}(u'-u)$ for
$u'\in P'$, and on $P'$
\[
 |\psi|\le W,\qquad\norm{s_\psi}\le G_V=\norm V_\triangle G .
\]
Consider the capped epigraph
\[
 \mathcal K=\{(u,s)\in\R^m\times\R:\ u\in P',\ \psi(u)\le s\le W+2\}.
\]
It is compact and convex. Let $a_0=(0,W+1)$ and
$r_{\mathcal K}=\min\{r_0,1/2\}$. A point $(u,s)$ with
$\norm{(u,s)-a_0}\le r_{\mathcal K}$ has $\norm u\le r_0$, hence $u\in P'$,
and $W+\tfrac12\le s\le W+\tfrac32$, hence $\psi(u)\le W<s\le W+2$. Thus the
ball of radius $r_{\mathcal K}$ about $a_0$ lies in $\mathcal K$, and
$\mathcal K$ is a convex body. Since $\norm u\le R_0$ and
$|s|\le W+2$ on $\mathcal K$, it lies in the ball of radius
$R_{\mathcal K}=R_0+W+2$ about the origin.

\emph{Separation.} Given rational $(u,s)$ and $\delta>0$, proceed as follows.
If $g_i^{\mathsf T}u>h_i$ for some $i$, return
$(g_i,0)/\norm{g_i}_\infty$; every point of $\mathcal K$ satisfies
$g_i^{\mathsf T}u'\le h_i<g_i^{\mathsf T}u$. Otherwise, if $s>W+2$, return
$(0,1)$. Otherwise $u\in P'$; compute $\psi(u)$ and $s_\psi(u)$. If $s<\psi(u)$,
return $(s_\psi(u),-1)$ divided by its maximum norm, which is at least one.
This is valid because for $(u',s')\in\mathcal K$ the subgradient inequality
gives
\[
 s'\ge\psi(u')\ge\psi(u)+s_\psi(u)^{\mathsf T}(u'-u)
 >s+s_\psi(u)^{\mathsf T}(u'-u).
\]
Otherwise $(u,s)\in\mathcal K$, and we assert $(u,s)\in S(\mathcal
K,\delta)$. Every returned vector separates $(u,s)$ from all of
$\mathcal K$, so the answers meet the weak separation requirement for every
$\delta>0$. The running time is polynomial in $\Delta+\langle u,s\rangle$.

\emph{Optimization.} Put
\[
 A_0=1+\frac{2W+1}{r_{\mathcal K}},\qquad
 C_0=A_0+1+G_V\Bigl(1+\frac{R_0}{r_0}\Bigr),\qquad
 \varepsilon=\min\Bigl\{\frac{r_{\mathcal K}}2,\frac\eta{C_0}\Bigr\}.
\]
Apply Lemma~\ref{lem:points-gls} with objective $c=(0,\ldots,0,-1)$ and
tolerance $\varepsilon$. The set $S(\mathcal K,-\varepsilon)$ contains the
ball of radius $r_{\mathcal K}-\varepsilon>0$ about $a_0$, so the algorithm
returns a rational $(u,s)\in S(\mathcal K,\varepsilon)$ with
$s\le s'+\varepsilon$ for every $(u',s')\in S(\mathcal K,-\varepsilon)$.

\emph{Lower bound.} Let $u_*$ minimize $\psi$ over $P'$, $\psi_*=\psi(u_*)$,
$z_*=(u_*,\psi_*)\in\mathcal K$, and $\lambda=\varepsilon/r_{\mathcal K}\le
\tfrac12$. For $\norm e\le\varepsilon$,
\[
 (1-\lambda)z_*+\lambda a_0+e=(1-\lambda)z_*+\lambda(a_0+e/\lambda),
\]
and $\norm{e/\lambda}\le r_{\mathcal K}$, so this point lies in
$\mathcal K$. Hence $(1-\lambda)z_*+\lambda a_0\in S(\mathcal
K,-\varepsilon)$. Its last coordinate is
$\psi_*+\lambda(W+1-\psi_*)\le\psi_*+\lambda(2W+1)$, because
$\psi_*\ge-W$. Therefore $s\le\psi_*+A_0\varepsilon$, and
$\ell=s-A_0\varepsilon$ satisfies $\ell\le\psi_*=\min_P\varphi$.

\emph{Feasible repair.} Some $(\bar u,\bar s)\in\mathcal K$ lies within
$\varepsilon$ of $(u,s)$. Put $\hat u=r_0u/(r_0+\varepsilon)$. For each row,
$g_i^{\mathsf T}u\le g_i^{\mathsf T}\bar u+\norm{g_i}\varepsilon\le
h_i+\norm{g_i}\varepsilon$, and $r_0\norm{g_i}\le r_0\norm{g_i}_1\le h_i$, so
\[
 g_i^{\mathsf T}\hat u\le\frac{r_0h_i+r_0\norm{g_i}\varepsilon}
 {r_0+\varepsilon}\le h_i .
\]
Thus $\hat u\in P'$. Moreover
$\hat u-\bar u=[r_0(u-\bar u)-\varepsilon\bar u]/(r_0+\varepsilon)$, and
$\norm{\bar u}\le R_0$, so
$\norm{\hat u-\bar u}\le\varepsilon(1+R_0/r_0)$. Both points lie in $P'$,
and the subgradient inequality at $\hat u$ gives
$\psi(\hat u)\le\psi(\bar u)+s_\psi(\hat u)^{\mathsf T}(\hat u-\bar u)\le
\psi(\bar u)+G_V\norm{\hat u-\bar u}$. Hence
\[
 \psi(\hat u)\le\psi(\bar u)+G_V\varepsilon(1+R_0/r_0)
 \le\bar s+G_V\varepsilon(1+R_0/r_0)
 \le s+\varepsilon+G_V\varepsilon(1+R_0/r_0).
\]
Return $\hat x=x_0+V\hat u\in P$ and $\ell$. Then
$\ell\le\min_P\varphi\le\varphi(\hat x)$ and
$\varphi(\hat x)-\ell\le C_0\varepsilon\le\eta$.

\emph{Complexity.} The data of Lemma~\ref{lem:points-relative}, and $A_0$,
$C_0$, $\varepsilon$, $R_{\mathcal K}$, have encoding length polynomial in
$\Delta+\langle\eta\rangle$. The separation routine runs in polynomial time
in its query length and $\Delta$. Lemma~\ref{lem:points-gls} and one final
evaluation give the stated bound.

(b) Lemma~\ref{lem:points-sparse} evaluates $\varphi$ and $\nabla\varphi$
exactly and supplies $W=\max\{1,V_f+\tau n\varrho^2\}$ and
$G=\max\{1,\norm T_\triangle G_f+2\tau n\varrho\}$. For a function that is
continuously differentiable near $P$ and convex on $P$, the gradient is a
subgradient on $P$: for $x,x'\in P$ the convex function
$\theta\mapsto\varphi(x+\theta(x'-x))$ on $[0,1]$ lies above its tangent at
$\theta=0$. Part (a) applies with $\Delta$ polynomial in $D$ and the stated
encoding lengths.
\end{proof}

\begin{remark}[An independently checkable lower bound]\label{rem:points-tangent}
If $\varphi$ is twice continuously differentiable with
$\norm{\nabla^2\varphi}\le H$ on $P$, a lower bound can also be certified
by linear programming duality. Let $\hat x\in P$ have
$\varphi(\hat x)-\min_P\varphi\le\delta$, let
$\gamma=\nabla\varphi(\hat x)$, and let $\hat\ell=\varphi(\hat x)+
\min_{x\in P}\gamma^{\mathsf T}(x-\hat x)$. Convexity gives
$\hat\ell\le\min_P\varphi$, and an optimal dual solution of the tangent
linear program certifies $\hat\ell$. If $x'$ attains the tangent minimum
and $\Delta_P\ge\max_{x,x'\in P}\norm{x-x'}$, the points
$\hat x+t(x'-\hat x)$, $0\le t\le1$, give
$t(\varphi(\hat x)-\hat\ell)\le\delta+Ht^2\Delta_P^2/2$. With
$K_0=\max\{1,H\Delta_P^2\}$, $\delta=\min\{\eta/2,\eta^2/(8K_0)\}$, and
$t=\min\{1,\eta/(2K_0)\}$ this yields $\varphi(\hat x)-\hat\ell\le\eta$.
A verifier checks feasibility of $\hat x$, the exact values $\varphi(\hat
x)$ and $\gamma$, and the dual certificate; it does not repeat the
ellipsoid computation. This certificate is optional and is not used in the
theorems.
\end{remark}

\paragraph{The selector lemma.}

\begin{proof}[Proof of Lemma~\ref{lem:selector}]
Write $\Phi=\varphi+\tau\norm{\cdot}^2$. Since $\varphi\ge\varphi_*$ on $P$,
every sublevel set of $\Phi$ in $P$ is bounded; it is closed by continuity.
Hence $\Phi$ attains its minimum on $P$, and the minimizer $x_\tau$ is unique
because $\tau\norm{\cdot}^2$ is strictly convex. Comparison with $p$ gives
\[
 \varphi(x_\tau)+\tau\norm{x_\tau}^2\le\varphi_*+\tau\norm p^2 ,
\]
so $\norm{x_\tau}\le\norm p\le R$ and
$\varphi(x_\tau)-\varphi_*\le\tau R^2\le1$; the last inequality holds
because $\varepsilon\le1\le R,\Gamma$ and $d\ge2$.

Let $s$ be the projection of $x_\tau$ onto the closed convex set $S$, and
$e=\norm{s-x_\tau}$. Comparison of $\Phi$ at $x_\tau$ and $s$ gives
\[
 \varphi(x_\tau)-\varphi_*\le\tau(\norm s^2-\norm{x_\tau}^2)
 \le\tau e(\norm s+\norm{x_\tau})\le\tau e(2\norm{x_\tau}+e).
\]
Since $p\in S$, $e\le\norm{p-x_\tau}\le2R$, and the right side is at most
$4R\tau e$. If every point of $S$ has norm at most $R$, then
$\norm s+\norm{x_\tau}\le2R$ and the bound is $2R\tau e$. The error bound
applies to $x_\tau$, so $e^d\le\Gamma^d\cdot4R\tau e$. If $e>0$, the choice
of $\tau$ gives
\[
 e^{d-1}\le4R\tau\Gamma^d=\Bigl(\frac{\varepsilon^2}{8R}\Bigr)^{d-1},
\]
so $e\le\varepsilon^2/(8R)$; this also holds if $e=0$, and in the bounded
variant with the factor $2$.

The point $p$ is the projection of the origin onto $S$, so
$p^{\mathsf T}(z-p)\ge0$ for all $z\in S$; in particular
$p^{\mathsf T}s\ge\norm p^2$. Using $\norm{x_\tau}\le\norm p$,
\[
 \norm{x_\tau-p}^2\le2\norm p^2-2p^{\mathsf T}x_\tau
 \le2p^{\mathsf T}(s-x_\tau)\le2Re\le\frac{\varepsilon^2}4 .
\]
Finally, for $x\in P$ and $0<\theta\le1$, the point
$x_\theta=x_\tau+\theta(x-x_\tau)$ lies in $P$, and convexity of $\varphi$
and the identity
$\norm{x_\theta}^2=(1-\theta)\norm{x_\tau}^2+\theta\norm x^2-
\theta(1-\theta)\norm{x-x_\tau}^2$ give
\[
 \Phi(x_\tau)\le\Phi(x_\theta)\le(1-\theta)\Phi(x_\tau)+\theta\Phi(x)
 -\tau\theta(1-\theta)\norm{x-x_\tau}^2 .
\]
Dividing by $\theta$ and letting $\theta\to0$ yields
$\Phi(x)-\Phi(x_\tau)\ge\tau\norm{x-x_\tau}^2$. If the left side is at most
$\tau\varepsilon^2/4$, then $\norm{x-x_\tau}\le\varepsilon/2$ and
$\norm{x-p}\le\varepsilon$.
\end{proof}

No differentiability of $\varphi$ is needed. The lemma does not assume
that $S$ is bounded; the factor $4$ pays for a projection $s$ that may lie
far from the origin.

\subsection{Proof of Theorem~\ref{thm:global-point}}\label{app:points-global}

Fix the data of Theorem~\ref{thm:global-point}: an explicit sparse
$f=\sum_\alpha f_\alpha x^\alpha$, convex on $\R^n$, and
$P=\{x:Ax\le b\}$. If $n=0$, the problem is the evaluation of a constant
and a check of the constraints, so assume $n\ge1$. Put
$D=\max\{2,\deg f\}$ and $\kappa_D=c_2(D)=(D+1)D^{D+2}$. Let
$\int\cdot\,da$ denote integration over $[0,1]^n$ with respect to Lebesgue
measure.

\begin{lemma}[Integrated quadratic minorant]\label{lem:points-minorant}
Let
\begin{align*}
 M&=\int\nabla^2f(a)\,da,\qquad
 g=\int\Bigl[\nabla f(a)-\frac2{\kappa_D}\nabla^2f(a)\,a\Bigr]da,\\
 m_0&=\int\Bigl[f(a)-\nabla f(a)^{\mathsf T}a+
 \frac1{\kappa_D}a^{\mathsf T}\nabla^2f(a)\,a\Bigr]da .
\end{align*}
\begin{enumerate}[label=(\alph*)]
\item $M$, $g$, and $m_0$ are rational, computable in time polynomial in
$L,D$, and $f(x)\ge m_0+g^{\mathsf T}x+x^{\mathsf T}Mx/\kappa_D$ for all
$x\in\R^n$.
\item $M\succeq0$, and $\ker M$ is the set of $d$ with $\nabla^2f(a)d=0$
for every $a\in\R^n$. For such $d$, $\nabla f(x)^{\mathsf T}d=
\nabla f(0)^{\mathsf T}d$ for all $x$.
\item Let $\Pi$ be the orthogonal projection onto the range of $M$
($\Pi=0$ if $M=0$) and $w=-(I-\Pi)\nabla f(0)$. Then $\Pi$ and $w$ are
rational and computable in polynomial time, and
\[
 f(x)=f(\Pi x)-w^{\mathsf T}x\quad(x\in\R^n),\qquad g=\Pi g-w .
\]
\item Let $Q_M$ be a positive integer with $Q_MM$ integral, let
$C_M\ge1$ bound the entries of $Q_MM$, and put
$\mu_0=1/(\kappa_DQ_M(nC_M)^{n-1})$. Then
$x^{\mathsf T}Mx/\kappa_D\ge\mu_0\norm{\Pi x}^2$ for all $x$.
\end{enumerate}
\end{lemma}

\begin{proof}
(a) Fix $a,x$ and let
$h(t)=f(a+t(x-a))-f(a)-t\nabla f(a)^{\mathsf T}(x-a)$. It has degree at
most $D$. By global convexity $h\ge0$ on $\R$, $h(0)=0$, and $h$ is convex,
so $0\le h(t)\le th(1)\le h(1)$ on $[0,1]$. Since
$h''(0)=(x-a)^{\mathsf T}\nabla^2f(a)(x-a)$,
Lemma~\ref{lem:points-interp}(c) gives
\[
 f(x)\ge f(a)+\nabla f(a)^{\mathsf T}(x-a)+
 \frac1{\kappa_D}(x-a)^{\mathsf T}\nabla^2f(a)(x-a).
\]
Expanding the right side in powers of $x$ and integrating over $a$ gives
the minorant. For computation, use the identities
\[
 (\nabla^2f(a)a)_i=\sum_\alpha f_\alpha\alpha_i(|\alpha|-1)a^{\alpha-e_i},
 \qquad
 a^{\mathsf T}\nabla^2f(a)a=\sum_\alpha f_\alpha|\alpha|(|\alpha|-1)a^\alpha,
\]
$\nabla f(a)^{\mathsf T}a=\sum_\alpha f_\alpha|\alpha|a^\alpha$, and
$\int a^\gamma\,da=\prod_i(\gamma_i+1)^{-1}$. Thus
\[
 M_{ij}=\sum_\alpha f_\alpha\alpha_i(\alpha_j-\delta_{ij})
 \!\int\! a^{\alpha-e_i-e_j}da,\quad
 g_i=\sum_\alpha f_\alpha\alpha_i\Bigl(1-\frac{2(|\alpha|-1)}{\kappa_D}\Bigr)
 \!\int\! a^{\alpha-e_i}da,
\]
and
$m_0=\sum_\alpha f_\alpha\bigl(1-|\alpha|+|\alpha|(|\alpha|-1)/\kappa_D\bigr)
\int a^\alpha da$, where terms with a negative exponent are omitted. Each
entry is a sum of at most one term per input monomial, with denominators
of $O(n\log(D+1)+D\log D)$ bits.

(b) $M$ is an integral of positive semidefinite matrices. If
$d^{\mathsf T}Md=0$, the continuous nonnegative polynomial
$a\mapsto d^{\mathsf T}\nabla^2f(a)d$ has zero integral over the cube, so it
vanishes on the cube and hence identically. Positive semidefiniteness gives
$\nabla^2f(a)d=0$ for all $a$. The converse is immediate. For such $d$, the
gradient of $x\mapsto\nabla f(x)^{\mathsf T}d$ is $\nabla^2f(x)d=0$.

(c) A rational basis $E$ of the range of $M$ gives
$\Pi=E(E^{\mathsf T}E)^{-1}E^{\mathsf T}$. Since $(I-\Pi)x\in\ker M$, part
(b) gives
\[
 f(x)-f(\Pi x)=\int_0^1\nabla f\bigl(\Pi x+\theta(I-\Pi)x\bigr)^{\mathsf T}
 (I-\Pi)x\,d\theta=\nabla f(0)^{\mathsf T}(I-\Pi)x=-w^{\mathsf T}x .
\]
Every $\nabla^2f(a)$ is symmetric and annihilates $\ker M$, so its range
lies in the range of $M$ and $(I-\Pi)\nabla^2f(a)=0$. By (b),
$(I-\Pi)\nabla f(a)=(I-\Pi)\nabla f(0)$. Applying $I-\Pi$ to the definition
of $g$ gives $(I-\Pi)g=-w$.

(d) This is Lemma~\ref{lem:points-integer-eigen}.
\end{proof}

The vector $w$ is the slope of $f$ along the common Hessian kernel. Kernel
directions with $w^{\mathsf T}d\ne0$ are directions of linear change, not
of invariance; part (c) records both.

\begin{lemma}[Boundedness, attainment, and a radius]\label{lem:points-radius}
Let $P\ne\emptyset$, let $a\in P$ be rational, and put $f_a=f(a)$. Consider
the linear system
\begin{equation}\label{eq:points-boundedness-lp}
 w=A^{\mathsf T}\lambda+Mz,\qquad\lambda\ge0 .
\end{equation}
\begin{enumerate}[label=(\alph*)]
\item If \eqref{eq:points-boundedness-lp} is infeasible, there is a
rational $d$ with $Ad\le0$, $Md=0$, and $w^{\mathsf T}d=1$, and
$f(a+td)=f_a-t$ for all $t\ge0$.
\item If it is feasible, let $(\lambda,z)$ be a rational solution and put
$v=Mz$, $\nu=\norm{\Pi g-v}_1$, and $\beta_0=m_0-\lambda^{\mathsf T}b$.
Then
\begin{equation}\label{eq:points-quotient}
 f(x)\ge\beta_0-\nu\norm{\Pi x}+\mu_0\norm{\Pi x}^2\qquad(x\in P).
\end{equation}
With $r_*=1+(\nu+|f_a-\beta_0|+1)/\mu_0$ and
$W_0=1+|\lambda^{\mathsf T}b|+\norm v_1r_*+|f(0)|+\norm{\nabla f(0)}_1r_*+
|f_a|$, every $x\in P$ with $f(x)\le f_a$ satisfies $\norm{\Pi x}\le r_*$
and $|w^{\mathsf T}x|\le W_0$.
\item In case (b), let $J=\binom{M}{w^{\mathsf T}}$, let $Q_J$ be a
positive integer with $Q_JJ$ integral, let $A_Z$ be obtained from $A$ by
multiplying each row by the least positive integer that makes it integral,
let $C\ge1$ bound the entries of $Q_JJ$ and $A_Z$, let
$m_1=\max_{ij}|M_{ij}|$, and put
\[
 R=1+\norm a_1+(nC)^{n-1}Q_J\bigl(\norm{Ja}_1+n^2m_1r_*+W_0\bigr).
\]
Every $y\in P$ with $f(y)\le f_a$ has a point $y'\in P$ with $f(y')=f(y)$
and $\norm{y'}<R$. Consequently $f$ attains its infimum on $P$, the optimal
set $S$ is nonempty, closed, and convex, and its minimum-norm point $p$
satisfies $\norm p<R$.
\item All quantities above are rational and computable in time polynomial
in $L,D$.
\end{enumerate}
\end{lemma}

\begin{proof}
(a) By Farkas' lemma, infeasibility gives $d$ with $Ad\le0$, $Md=0$, and
$w^{\mathsf T}d>0$; a linear program finds one with $w^{\mathsf T}d=1$. Then
$a+td\in P$ for $t\ge0$, $\Pi d=0$, and Lemma~\ref{lem:points-minorant}(c)
gives $f(a+td)=f(\Pi a)-w^{\mathsf T}a-t=f_a-t$.

(b) For $x\in P$, $\lambda\ge0$ and $Ax\le b$ give
$w^{\mathsf T}x=\lambda^{\mathsf T}Ax+z^{\mathsf T}Mx\le
\lambda^{\mathsf T}b+v^{\mathsf T}\Pi x$, because $v$ lies in the range of
$M$. By Lemma~\ref{lem:points-minorant},
\[
 f(x)\ge m_0+(\Pi g)^{\mathsf T}\Pi x-w^{\mathsf T}x+\mu_0\norm{\Pi x}^2
 \ge\beta_0+(\Pi g-v)^{\mathsf T}\Pi x+\mu_0\norm{\Pi x}^2 ,
\]
and $|(\Pi g-v)^{\mathsf T}\Pi x|\le\nu\norm{\Pi x}$. This proves
\eqref{eq:points-quotient}. If $f(x)\le f_a$ and $\rho=\norm{\Pi x}>r_*$,
then $\rho>1$ and $\mu_0\rho-\nu>|f_a-\beta_0|+1$, so
$\mu_0\rho^2-\nu\rho>|f_a-\beta_0|+1>f_a-\beta_0$, contradicting
\eqref{eq:points-quotient}. Hence $\norm{\Pi x}\le r_*$. The upper bound
$w^{\mathsf T}x\le\lambda^{\mathsf T}b+\norm v_1r_*$ was shown above. For
the lower bound, global convexity at $0$ and
Lemma~\ref{lem:points-minorant}(c) give
\[
 w^{\mathsf T}x=f(\Pi x)-f(x)\ge f(0)+\nabla f(0)^{\mathsf T}\Pi x-f_a
 \ge f(0)-\norm{\nabla f(0)}_1r_*-f_a .
\]
The constraints $Ax\le b$ alone give only the upper bound.

(c) Let $y\in P$ with $f(y)\le f_a$. Since $My=M\Pi y$,
$\norm{My}_1\le nm_1\norm{\Pi y}_1\le n^2m_1r_*$, and
$\norm{Q_JJy}_1\le Q_J(n^2m_1r_*+W_0)$. The slice
$Z_y=\{x\in P:Q_JJx=Q_JJy\}$ contains $y$. Apply
Lemma~\ref{lem:integer-hoffman} with inequality matrix $A_Z$ and equality
matrix $Q_JJ$, both integral, to the point $a\in P$. The projection $y'$ of
$a$ onto $Z_y$ satisfies
\[
 \norm{y'}\le\norm a+(nC)^{n-1}\norm{Q_JJ(a-y)}<R .
\]
From $My'=My$ we get $y'-y\in\ker M$ and $\Pi y'=\Pi y$; with
$w^{\mathsf T}y'=w^{\mathsf T}y$, Lemma~\ref{lem:points-minorant}(c) gives
$f(y')=f(y)$.

By (b), $f$ is bounded below on $P$. If $a$ is optimal, its own
representative is $a$, of norm at most $\norm a_1<R$. Otherwise a minimizing
sequence eventually lies in $\{f\le f_a\}$; replace it by representatives
in the compact set $P\cap\{\norm x\le R\}$, extract a convergent
subsequence, and use continuity. Either way $f$ attains its infimum at a
point of norm less than $R$. The optimal set is a sublevel set of a convex
continuous function on a closed convex set, so its minimum-norm point exists
and has norm less than $R$. The argument does not assume attainment in
order to construct $R$.

(d) The point $a$, the system \eqref{eq:points-boundedness-lp}, and its
solution or Farkas certificate come from linear programs of polynomial size.
Lemma~\ref{lem:points-minorant} supplies $M,g,m_0,\Pi,w,\mu_0$. The values
$f_a$, $f(0)$, $\nabla f(0)$ and all products in $r_*,W_0,R$ have
polynomial encoding length; the power $(nC)^{n-1}$ has
$(n-1)\log_2(nC)$ bits.
\end{proof}

\begin{lemma}[Global error modulus]\label{lem:points-error}
Assume that $S\ne\emptyset$ and that its minimum-norm point $p$ satisfies
$\norm p\le R$ for a rational $R\ge1$. Let
$\mathcal E=\{\alpha-e_i:f_\alpha\ne0,\ \alpha_i>0\}$ and define
$N\in\Q^{\mathcal E\times n}$ by
\[
 N_{\gamma,i}=(\gamma_i+1)f_{\gamma+e_i},
\]
where $f_{\gamma+e_i}=0$ if $\gamma+e_i$ is not an exponent of $f$. Let
$Q_N$ be a positive integer with $Q_NN$ integral, let $C_N\ge1$ bound the
entries of $Q_NN$ and $A_Z$, let $k_N=\max\{1,|\mathcal E|\}$, and put
\[
 V_f=\sum_\alpha|f_\alpha|(R+2)^{|\alpha|},\quad
 G_f=\sum_{|\alpha|\ge1}|f_\alpha||\alpha|(R+2)^{|\alpha|-1},\quad
 h_f=\max\{1,2V_f+2G_f(R+1)\},
\]
\[
 C_*=1+c_1(D)\bigl(h_f+3^Dc_0(D)\bigr),\qquad
 \Gamma=\max\bigl\{1,(nC_N)^{n-1}Q_Nk_N\beta_{D-1}^nC_*\bigr\}.
\]
Then $S=\{x\in P:Nx=Np\}$ and
\[
 \dist(x,S)\le\Gamma\bigl(f(x)-f_*\bigr)^{1/D}\qquad
 (x\in P,\ f(x)-f_*\le1).
\]
The number $\Gamma$ is computable in time polynomial in $L,D,\langle R\rangle$.
\end{lemma}

\begin{proof}
\emph{Coefficient rows.} For every $d$ and $z$,
\[
 \nabla f(z)^{\mathsf T}d=\sum_i d_i\sum_{\alpha_i>0}\alpha_if_\alpha
 z^{\alpha-e_i}=\sum_{\gamma\in\mathcal E}(Nd)_\gamma z^\gamma ,
\]
and distinct $\gamma$ give distinct monomials. Hence $Nd=0$ if and only if
$\nabla f(z)^{\mathsf T}d=0$ for all $z$, that is, if and only if
$f(z+\theta d)=f(z)$ for all $z$ and all real $\theta$.

\emph{A Bregman estimate.} Fix $x\in P$ with
$E=f(x)-f_*\in[0,1]$, put $d=x-p$, and let
\[
 h(u)=f(p+u)-f(p)-\nabla f(p)^{\mathsf T}u .
\]
Global convexity makes $h$ convex and nonnegative on $\R^n$. Optimality of
$p$ over $P$ gives $\nabla f(p)^{\mathsf T}d\ge0$, and convexity gives
$\nabla f(p)^{\mathsf T}d\le E$. For $0\le\theta\le1$,
\[
 0\le h(\theta d)\le f(p+\theta d)-f_*\le\theta E\le E .
\]
The polynomial $\theta\mapsto h(\theta d)$ has degree at most $D$, so
Lemma~\ref{lem:points-interp}(a) gives
$h(2\theta d)\le c_0(D)(1+2|\theta|)^DE$ for every real $\theta$.

Let $c\in[0,1]^n$ and $\pi_c(\theta)=h(c-p+\theta d)$. Since
$c-p+\theta d$ is the midpoint of $2(c-p)$ and $2\theta d$, convexity of $h$
gives
\begin{equation}\label{eq:points-midpoint}
 0\le\pi_c(\theta)\le\tfrac12h(2(c-p))+\tfrac12h(2\theta d).
\end{equation}
The points $p$ and $2c-p$ lie in the box of radius $R+2$, and
$\norm{c-p}_\infty\le R+1$. For $\norm y_\infty\le R+2$ every monomial
satisfies $|y^\alpha|\le(R+2)^{|\alpha|}$ and
$\sum_i|\partial_iy^\alpha|\le|\alpha|(R+2)^{|\alpha|-1}$, so $|f(y)|\le V_f$
and $\norm{\nabla f(y)}_1\le G_f$. Hence
\[
 h(2(c-p))\le|f(2c-p)|+|f(p)|+2\norm{\nabla f(p)}_1\norm{c-p}_\infty\le h_f .
\]
Only the norm of $p$ enters here; the point $x$ may be arbitrarily far
away.

If $E>0$, then on $0\le\theta\le E^{-1/D}$ we have
$E(1+2\theta)^D\le E(3E^{-1/D})^D=3^D$, so
$0\le\pi_c\le h_f+3^Dc_0(D)$ there. Lemma~\ref{lem:points-interp}(b) with
$\varrho=E^{-1/D}$ gives
$|\pi_c'(0)|\le c_1(D)(h_f+3^Dc_0(D))E^{1/D}$. Since
$\pi_c'(0)=(\nabla f(c)-\nabla f(p))^{\mathsf T}d$ and
$E\le E^{1/D}$,
\begin{equation}\label{eq:points-gradient-residual}
 |\nabla f(c)^{\mathsf T}d|\le C_*E^{1/D}\qquad(c\in[0,1]^n).
\end{equation}
If $E=0$, then $h(\theta d)=0$ on $[0,1]$, so this polynomial vanishes
identically, and \eqref{eq:points-midpoint} bounds $\pi_c$ on the whole real
line. A bounded polynomial is constant, so $\pi_c'(0)=0$; also
$\nabla f(p)^{\mathsf T}d=0$. Thus \eqref{eq:points-gradient-residual}
holds with right side zero.

\emph{From residuals to coefficients.} The polynomial
$z\mapsto\nabla f(z)^{\mathsf T}d=\sum_\gamma(Nd)_\gamma z^\gamma$ has
degree at most $D-1\ge1$ in each variable, and
\eqref{eq:points-gradient-residual} bounds it on $[0,1]^n$.
Lemma~\ref{lem:points-interp}(d) gives
\[
 |(N(x-p))_\gamma|\le\beta_{D-1}^nC_*E^{1/D}\qquad(\gamma\in\mathcal E).
\]

\emph{The optimal slice.} If $x\in S$, then $E=0$ and $N(x-p)=0$.
Conversely, if $x\in P$ and $N(x-p)=0$, then $f(x)=f(p)$ by invariance, so
$x\in S$. Hence $S=\{x\in P:Q_NNx=Q_NNp\}$. Apply
Lemma~\ref{lem:integer-hoffman} with the integer matrices $A_Z$ and $Q_NN$:
\[
 \dist(x,S)\le(nC_N)^{n-1}\norm{Q_NN(x-p)}
 \le(nC_N)^{n-1}Q_N\sqrt{k_N}\,\beta_{D-1}^nC_*E^{1/D}\le\Gamma E^{1/D}.
\]
If $\mathcal E$ is empty, $f$ is constant, $S=P$, and the bound is trivial.

\emph{Complexity.} The matrix $N$ has at most $n$ times as many rows as $f$
has monomials, and its entries are products of input coefficients and
exponents. The numbers $V_f,G_f,h_f$ are sums of at most one term per
monomial, each of $O(\langle f_\alpha\rangle+D\langle R\rangle)$ bits. The
constants $c_0(D),c_1(D),3^D$ have $O(D\log(D+1))$ bits and
$\beta_{D-1}^n$ has $O(nD\log(D+1))$ bits.
\end{proof}

\begin{proof}[Proof of Theorem~\ref{thm:global-point}]
\emph{Part (a).} Decide emptiness of $P$ by linear programming. If
$P\ne\emptyset$, compute a rational $a\in P$ and the data of
Lemma~\ref{lem:points-minorant}, and solve
\eqref{eq:points-boundedness-lp}. By Lemma~\ref{lem:points-radius}, the
system is infeasible exactly when $f$ is unbounded below on $P$, and in that
case the lemma supplies the certificate $(a,d)$.

\emph{Part (b).} If \eqref{eq:points-boundedness-lp} is feasible,
Lemma~\ref{lem:points-radius}(c) gives attainment and a rational $R$ with
$\norm p<R$, and Lemma~\ref{lem:points-error} with this $R$ gives $\Gamma$.
Both are computed in time polynomial in $L,D$, so $\log R$ and $\log\Gamma$
are polynomial in $L,D$.

\emph{Part (c).} Let $\varepsilon=2^{-q}$ and define $\tau$ by
Lemma~\ref{lem:selector} with $d=D$ and the computed $R,\Gamma$. Then
\[
 \log_2(1/\tau)\le(2D-2)q+(3D-1)+D\log_2R+D\log_2\Gamma ,
\]
which is polynomial in $L,D,q$. By Lemma~\ref{lem:selector},
$\norm{x_\tau}\le\norm p<R$, so $x_\tau$ lies in the bounded polytope
$Q=P\cap[-R,R]^n$ and minimizes $f+\tau\norm{\cdot}^2$ over $Q$ as well.
Apply Lemma~\ref{lem:convex-value} to $Q$, $\varrho=R$, the function
$\varphi=f+\tau\norm{\cdot}^2$, and $\eta=\tau\varepsilon^2/4$. The
evaluation and the bounds $W,G$ come from Lemma~\ref{lem:points-sparse} with
$T=I$, $t=0$; the composition $f(x_0+Vu)$ inside
Lemma~\ref{lem:convex-value} is evaluated, not expanded. The returned
rational $x_q\in Q\subseteq P$ satisfies
$\varphi(x_q)-\varphi(x_\tau)\le\eta$, hence $\norm{x_q-p}\le\varepsilon$.

For a value enclosure of width $\delta=2^{-q_f}$, let
$G_Q=\max\{1,\sum_{|\alpha|\ge1}|f_\alpha||\alpha|R^{|\alpha|-1}\}$, which
bounds $\norm{\nabla f}$ on $[-R,R]^n$. Run the preceding procedure with
accuracy $\varepsilon'=\min\{2^{-q},\delta/(2G_Q)\}$. The returned point
$x_q$ is within $2^{-q}$ of $p$, and since the segment $[p,x_q]$ lies in $Q$,
$f(x_q)-f_*\le G_Q\varepsilon'\le\delta/2$. Apply
Lemma~\ref{lem:convex-value} to $f$ on $Q$ with $\eta=\delta/2$ to get
$\ell\le\min_Qf$. Because $p\in Q$, $\min_Qf=f_*$. Thus
$\ell\le f_*\le f(x_q)\le\ell+\delta$. The encoding lengths involved are
polynomial in $L,D,q+q_f$.
\end{proof}

\begin{proof}[Proof of Corollary~\ref{cor:points-bounded}]
For bounded $P\subseteq[-\varrho,\varrho]^n$, the function $f$ attains its
minimum by compactness, and $R_x=\max\{1,n\varrho\}$ bounds the norm of every
point of $P$. Lemma~\ref{lem:points-error} applies with $R=R_x$, without
Lemmas~\ref{lem:points-minorant} and~\ref{lem:points-radius}. Every optimal
point has norm at most $R_x$, so Lemma~\ref{lem:selector} holds with the
factor $2$, and the proof of part (c) above applies with $Q=P$. For fixed
$D$, every bound polynomial in $L,D,q$ is polynomial in $L+q$.
\end{proof}

\begin{proof}[Proof of the claim in Remark~\ref{rem:points-expansion}]
Let $n=2r\ge4$, $f(y)=\sum_iy_i^{2r}$, and $U=I-\tfrac2n\mathbf1\mathbf1^{\mathsf
T}$. Then $U^{\mathsf T}U=I-\tfrac4n\mathbf1\mathbf1^{\mathsf T}+
\tfrac4{n^2}n\mathbf1\mathbf1^{\mathsf T}=I$, and every entry of $U$ is
nonzero because $n\ge4$. For an exponent $\beta$ with $|\beta|=r$, the
coefficient of $y^{2\beta}$ in $f(Uy)=\sum_i(\sum_jU_{ij}y_j)^{2r}$ is
\[
 \frac{(2r)!}{\prod_j(2\beta_j)!}\sum_i\prod_jU_{ij}^{2\beta_j}>0 ,
\]
a sum of positive terms. There are $\binom{3r-1}r\ge2^r$ such exponents,
while $f$ has $n$ monomials of degree $2r$.
\end{proof}

\subsection{Proof of Theorem~\ref{thm:cubic-point}}\label{app:points-cubic}

\begin{lemma}[Reflection]\label{lem:points-reflection}
Let $H:\R^m\to\R^{m\times m}$ be an affine map into symmetric matrices, let
$\Omega\subseteq\R^m$ be convex with $H(u)\succeq0$ for $u\in\Omega$, and let
$c\in\Omega$ and $\theta\in(0,1]$ be such that $c-\theta(u-c)\in\Omega$ for
every $u\in\Omega$. Then $0\preceq H(u)\preceq(1+\theta^{-1})H(c)$ for
$u\in\Omega$; in particular $\ker H(c)\subseteq\ker H(u)$. If $\Omega$
contains the ball of radius $r$ about $c$ and lies in the ball of radius
$R\ge r$ about $c$, one can take $\theta=r/R$. If $\Omega$ is a box and $c$
is its center, one can take $\theta=1$.
\end{lemma}

\begin{proof}
With $u'=c-\theta(u-c)\in\Omega$ we have
$c=\frac\theta{1+\theta}u+\frac1{1+\theta}u'$, so affinity gives
$H(c)=\frac\theta{1+\theta}H(u)+\frac1{1+\theta}H(u')\succeq
\frac\theta{1+\theta}H(u)$. If $H(c)d=0$, then $0\le d^{\mathsf T}H(u)d\le0$
and positive semidefiniteness gives $H(u)d=0$. In the ball case,
$\norm{u'-c}=\theta\norm{u-c}\le r$; in the box case $u'$ is the reflection
of $u$ through the center.
\end{proof}

For a polynomial $\varphi$ of degree at most three write
$H(u)=\nabla^2\varphi(u)$; it is affine in $u$. Let $\mathsf D^3\varphi$ be
the constant, fully symmetric third derivative, so that
$H(u)-H(v)=\mathsf D^3\varphi[u-v]$ and Taylor's formula is exact:
\begin{equation}\label{eq:points-cubic-taylor}
 \varphi(v+d)=\varphi(v)+\nabla\varphi(v)^{\mathsf T}d+\tfrac12d^{\mathsf
 T}H(v)d+\tfrac16\mathsf D^3\varphi[d,d,d].
\end{equation}

\begin{proof}[Proof of Lemma~\ref{lem:points-cubic-transverse}]
Put $a=d^{\mathsf T}H(v)d$ and $b=d^{\mathsf T}H(u)d$, both nonnegative.
Since $b=a+\mathsf D^3\varphi[d,d,d]$, \eqref{eq:points-cubic-taylor} gives
$E=\nabla\varphi(v)^{\mathsf T}d+(2a+b)/6$. First-order optimality of $v$
over the convex set $\Omega$ gives $\nabla\varphi(v)^{\mathsf T}d\ge0$, so
$\sigma=a+b\le6E$. Symmetry of $\mathsf D^3\varphi$ gives
\[
 d^{\mathsf T}H(c)d=a+\mathsf D^3\varphi[c-v,d,d]
 =a+(c-v)^{\mathsf T}\bigl(H(u)-H(v)\bigr)d .
\]
For a positive semidefinite $Q$, $\norm{Qd}^2=d^{\mathsf T}Q^2d\le
\norm Q\,d^{\mathsf T}Qd$. Hence $\norm{H(u)d}+\norm{H(v)d}\le
\sqrt{\Lambda b}+\sqrt{\Lambda a}\le\sqrt{2\Lambda\sigma}$, and
$|d^{\mathsf T}H(c)d|\le\sigma+\rho_1\sqrt{2\Lambda\sigma}$. Also
$\sigma\le2\Lambda\rho_2^2$, so $\sigma\le\rho_2\sqrt{2\Lambda\sigma}$.
Therefore $|d^{\mathsf T}H(c)d|\le(\rho_1+\rho_2)\sqrt{2\Lambda\sigma}$, and
squaring with $\sigma\le6E$ proves the lemma.
\end{proof}

\begin{proof}[Proof of Lemma~\ref{lem:points-cubic-slice}]
Let $\psi(u)=\varphi(u)-g_c^{\mathsf T}u$. For $u\in\Omega$ the segment from
$c$ to $u$ lies in $\Omega$, and
$\nabla\psi(u)=\int_0^1H(c+\theta(u-c))(u-c)\,d\theta$. Each $H$ on $\Omega$
is symmetric and annihilates $K$, so $\nabla\psi(u)\in K^\perp$.

Let $w$ minimize $\varphi$ over $\Omega$, put $d=w-v$, $a=d^{\mathsf
T}H(v)d$, and $b=d^{\mathsf T}H(w)d$. Since $\varphi(w)=\varphi(v)$, the
computation in the previous proof gives
$0=\nabla\varphi(v)^{\mathsf T}d+(2a+b)/6$ with three nonnegative terms, so
$a=b=0$ and $H(v)d=H(w)d=0$. The identity in the previous proof gives
$d^{\mathsf T}H(c)d=0$, hence $H(c)d=0$ and $d\in K$. Then
$\psi(w)-\psi(v)=\int_0^1\nabla\psi(v+\theta d)^{\mathsf T}d\,d\theta=0$
and $0=\varphi(w)-\varphi(v)=g_c^{\mathsf T}d$. Conversely, if $w\in\Omega$,
$H(c)(w-v)=0$, and $g_c^{\mathsf T}(w-v)=0$, then $w-v\in K$,
$\psi(w)=\psi(v)$, and $\varphi(w)=\varphi(v)$.

For the estimate, let $u\in\Omega$ and $d=u-v$. Since
$\nabla\psi\in K^\perp$, $\nabla\psi^{\mathsf T}d=\nabla\psi^{\mathsf
T}\Pi_Kd$, and the mean value bound on $[c,z]\subseteq\Omega$ gives
$\norm{\nabla\psi(z)}=\norm{\nabla\varphi(z)-\nabla\varphi(c)}\le\Lambda
R_\Omega$. Hence
\[
 |g_c^{\mathsf T}d|=\Bigl|\varphi(u)-\varphi(v)-\int_0^1\nabla\psi(v+\theta
 d)^{\mathsf T}\Pi_Kd\,d\theta\Bigr|\le\varphi(u)-\varphi(v)+\Lambda
 R_\Omega\norm{\Pi_Kd}. \qedhere
\]
\end{proof}

\begin{proof}[Proof of Theorem~\ref{thm:cubic-point}]
\emph{Reduction.} Compute $\varrho=\max\{1,\max_j\max_{x\in P}|x_j|\}$ by
$2n$ linear programs and apply Lemma~\ref{lem:points-relative}. If $m=0$,
then $P=S=\{x_0\}$; take $\Gamma_P=1$ and output $x_0$. Otherwise let
$\varphi(u)=f(x_0+Vu)$, a cubic in $u$, and $P'=\{Gu\le h\}$ with inner
radius $r_0$ and outer bound $R_0$ about $0$. Convexity of $f$ on $P$ makes
$\varphi$ convex on the full-dimensional set $P'$, so
$H(u)=\nabla^2\varphi(u)=V^{\mathsf T}\nabla^2f(x_0+Vu)V\succeq0$ on the
interior of $P'$ and, by continuity, on $P'$. This step needs the reduction:
convexity on a lower-dimensional polytope does not make the ambient Hessian
positive semidefinite. We never expand $\varphi$; only
$H_c=H(0)=V^{\mathsf T}\nabla^2f(x_0)V$, $g_c=\nabla\varphi(0)=V^{\mathsf
T}\nabla f(x_0)$, and evaluations of $f$ are used.

\emph{Center Hessian.} Lemma~\ref{lem:points-reflection} with $c=0$ and
$\theta=r_0/R_0$ gives $0\preceq H(u)\preceq\omega H_c$ on $P'$, where
$\omega=1+R_0/r_0$, and $K=\ker H_c$ lies in every $\ker H(u)$, $u\in P'$.
Put $\Lambda=\max\{1,\omega\max_i\sum_j|(H_c)_{ij}|\}$; then
$\norm{H(u)}\le\Lambda$ on $P'$. Let $\Delta$ be a positive integer that
clears the denominators of $H_c$, $g_c$, and $G$, let
$C_*\ge1$ bound the entries of $T_Z=\Delta\binom{H_c}{g_c^{\mathsf T}}$ and
$\Delta G$, and let $S'$ be the optimal set of $\varphi$ on $P'$, so that
$S=x_0+VS'$ and $f(x_0+Vu)-f_*=\varphi(u)-\varphi_*$.

\emph{Affine branch.} If $H_c=0$, then $H$ vanishes on the open set
$\operatorname{int}P'$, hence everywhere, and $\varphi$ is affine with
gradient $g_c$. If $g_c=0$, then $S'=P'$ and $\Gamma_P=1$ works. Otherwise
$S'=\{u\in P':g_c^{\mathsf T}u=g_c^{\mathsf T}v\}$ for any $v\in S'$, and
$\Delta g_c^{\mathsf T}(u-v)=\Delta E$ with $E=\varphi(u)-\varphi_*$.
Lemma~\ref{lem:integer-hoffman} gives $\dist(u,S')\le(mC_*)^{m-1}\Delta E$,
which is at most $(mC_*)^{m-1}\Delta E^{1/4}$ for $E\le1$.

\emph{Curved branch.} Let $H_c\ne0$ have rank $r$. The product of its
nonzero eigenvalues is the sum of its principal minors of order $r$, a
positive rational with denominator dividing $\Delta^r$; so it is at least
$\Delta^{-r}$. Each eigenvalue is at most $\norm{H_c}\le\Lambda$. Hence the
least positive eigenvalue $\lambda$ satisfies
$\lambda\ge\lambda_0=\Delta^{-m}\Lambda^{-(m-1)}$. Let $v\in S'$, $u\in P'$,
$d=u-v$, and $E=\varphi(u)-\varphi(v)\le1$. With $\rho_1=R_0\ge\norm v$
and $\rho_2=2R_0\ge\norm d$, Lemma~\ref{lem:points-cubic-transverse} gives
\[
 \lambda_0^2\norm{\Pi_Kd}^4\le(d^{\mathsf T}H_cd)^2\le108\Lambda R_0^2E,
\]
so $\norm{\Pi_Kd}\le\rho_TE^{1/4}$ with
$\rho_T=\max\{1,108\Lambda R_0^2/\lambda_0^2\}$. By
Lemma~\ref{lem:points-cubic-slice} with $\Omega=P'$, $c=0$, and
$R_\Omega=R_0$, the optimal set is $S'=\{u\in P':T_Zu=T_Zv\}$, and
$|g_c^{\mathsf T}d|\le E+\Lambda R_0\norm{\Pi_Kd}$. Since
$H_cd=H_c\Pi_Kd$,
\[
 \norm{T_Zd}\le\Delta\bigl(\Lambda\norm{\Pi_Kd}+E+\Lambda
 R_0\norm{\Pi_Kd}\bigr)\le\Delta\bigl(1+\Lambda(1+R_0)\rho_T\bigr)E^{1/4}.
\]
Lemma~\ref{lem:integer-hoffman}, applied to $P'$ with the integer matrices
$\Delta G$ and $T_Z$, gives $\dist(u,S')\le\Gamma_uE^{1/4}$ with
\[
 \Gamma_u=(mC_*)^{m-1}\Delta\bigl(1+\Lambda(1+R_0)\rho_T\bigr).
\]

\emph{Original coordinates.} For $x=x_0+Vu\in P$,
$\dist(x,S)\le\norm V\dist(u,S')$. Hence
$\Gamma_P=\max\{1,\norm V_\triangle\}\Gamma_u$ (or the affine-branch
constant multiplied by the same factor) satisfies the error bound. All
quantities are rational, computed in polynomial time, and have polynomial
encoding length.

\emph{Output.} For a set approximation, apply Lemma~\ref{lem:convex-value} to
$f$ on $P$ with $\eta=(2^{-q}/\Gamma_P)^4\le1$; the returned $\hat x$ has
$\dist(\hat x,S)\le\Gamma_P\eta^{1/4}=2^{-q}$. For the fixed selector, let
$R_x=\max\{1,n\varrho\}$, which bounds the norm of every point of $P$, and
apply Lemma~\ref{lem:selector} with $d=4$, $\Gamma=\Gamma_P$, $R=R_x$, and the
factor $2$:
\[
 \tau=\frac{\varepsilon^6}{1024R_x^4\Gamma_P^4},\qquad
 \eta=\frac{\tau\varepsilon^2}4,\qquad\varepsilon=2^{-q}.
\]
Lemma~\ref{lem:convex-value}, applied to $f+\tau\norm x^2$ on $P$ in the
original coordinates, returns $\hat x\in P$ with regularized gap at most
$\eta$, hence $\norm{\hat x-p}\le\varepsilon$. The penalty is the original
norm $\norm x^2=\norm{x_0+Vu}^2$, not $\norm u^2$; the latter would select
a different optimizer in general. The encoding lengths of $\tau$ and $\eta$
are $O(q)$ plus a polynomial in $L$.
\end{proof}

\subsection{Boundary constructions}\label{app:points-boundary}

\paragraph{Treewidth.}
Every objective below for which a treewidth bound is claimed is a sum of
terms in at most three variables each. Every monomial of the expanded sum
occurs in some term, so every edge of the interaction graph joins two
variables of one term. Hence a tree of bags in which the variables of each
term lie together in some bag, and each variable occurs in a connected set of
bags, is a tree decomposition of the interaction graph; bags of at most three
variables give treewidth at most two.

\paragraph{Two generic devices.}
The reductions combine a strongly convex base problem with a one-sided
sign test and a quartic amplifier.

\begin{lemma}[Sign test]\label{lem:points-sign-test}
Let $Y\subseteq\R^k$ be a box and $\Phi$ a polynomial with
$\nabla^2\Phi\succeq\kappa I$ on $Y$ for some $\kappa\ge1/8$, and let $y^0$ be
its minimizer over $Y$. Let $\zeta(y)=\sigma^{\mathsf T}y+\sigma_0$ be affine
with $\norm\sigma\le1$ and $|\zeta|\le2$ on $Y$. Put $\epsilon_0=1/32$ and
\[
 H(y,\theta)=\Phi(y)+\theta^2+\epsilon_0\theta\zeta(y)\qquad
 \text{on }Y\times[0,1].
\]
Then $\nabla^2H\succeq\tfrac1{16}I$ on $Y\times[0,1]$, so $H$ has a unique
minimizer $(y^*,\theta^*)$. If $\zeta(y^0)\ge0$, then $(y^*,\theta^*)=(y^0,0)$;
if $\zeta(y^0)<0$, then $0<\theta^*<1$.
\end{lemma}

\begin{proof}
The Hessian of $\Phi(y)+\theta^2$ is at least $\min\{\kappa,2\}I\succeq
\tfrac18I$, and the Hessian of $\epsilon_0\theta\zeta(y)$ has spectral norm
$\epsilon_0\norm\sigma\le\tfrac1{32}$. On the face $\theta=0$, $H=\Phi$, with
the unique minimizer $y^0$. At $(y^0,0)$ the partial derivatives in $y$ equal
$\nabla\Phi(y^0)$, which satisfies the first-order condition for $\Phi$ on
$Y$, and $\partial_\theta H=\epsilon_0\zeta(y^0)$. If $\zeta(y^0)\ge0$, the
point satisfies the first-order condition on $Y\times[0,1]$ and is the
minimizer by convexity. If $\zeta(y^0)<0$, the direction $+e_\theta$ is a
feasible descent direction at $(y^0,0)$, and no other point of the face
$\theta=0$ can be optimal, since its restriction has only the minimizer
$y^0$; so $\theta^*>0$. At $\theta=1$,
$\partial_\theta H\ge2-2\epsilon_0>0$, which excludes $\theta^*=1$.
\end{proof}

\begin{lemma}[Quartic amplifiers]\label{lem:points-amplifier}
Let $\eta_0=1/192$.
\begin{enumerate}[label=(\alph*)]
\item Let $H$ be a polynomial on a box $Z$ with $\nabla^2H\succeq\frac1{16}I$
on $Z$, unique minimizer $z^*$, and a distinguished coordinate
$\theta\in[0,1]$. Then $F'(z,y)=H(z)+\eta_0\theta^2(y-1)^2$ is convex on
$Z\times[0,1]$, and its optimal set is $\{z^*\}\times[0,1]$ if $\theta^*=0$
and $\{(z^*,1)\}$ if $\theta^*>0$.
\item Let $H_\pm$ be polynomials on boxes $Z_\pm$ with Hessians at least
$\frac1{16}I$, unique minimizers $z_\pm^*$, and distinguished coordinates
$\theta_\pm\in[0,1]$, such that exactly one of $\theta_+^*,\theta_-^*$ is
positive. Then
\[
 F(z_+,z_-,y)=H_+(z_+)+H_-(z_-)+\eta_0\bigl[\theta_+^2(y-1)^2+
 \theta_-^2y^2\bigr]
\]
is convex on $Z_+\times Z_-\times[0,1]$ and has a unique minimizer, whose
$y$-coordinate is $1$ if $\theta_+^*>0$ and $0$ if $\theta_-^*>0$.
\end{enumerate}
\end{lemma}

\begin{proof}
For real $\theta,y,\gamma$ and a direction $(\pi,\varsigma)$ in
$(\theta,y)$,
\begin{equation}\label{eq:points-amplifier-identity}
 \frac{d^2}{dt^2}\Bigl[(\theta+t\pi)^2(y+t\varsigma-\gamma)^2\Bigr]_{t=0}
 =2\bigl(\theta\varsigma+2(y-\gamma)\pi\bigr)^2-6(y-\gamma)^2\pi^2 ,
\end{equation}
as direct expansion shows. For $y\in[0,1]$ and $\gamma\in\{0,1\}$ the
negative term is at least $-6\pi^2$. Multiplied by $\eta_0$, it is at least
$-\frac1{32}\pi^2$, which the base Hessian absorbs; in (b) the two losses
affect the distinct coordinates $\theta_+$ and $\theta_-$. Hence the Hessian
quadratic form is at least $\frac1{32}$ times the squared norm of the base
part of the direction, with no division by $\theta$. This proves convexity,
including where the amplitudes vanish.

(a) Since $F'\ge H\ge H(z^*)$, a point is optimal exactly when $z=z^*$ and
$\theta^{*2}(y-1)^2=0$, and $y=1$ always attains the bound.

(b) $F$ is at least the sum of the two base minima. If $\theta_+^*>0$ and
$\theta_-^*=0$, the point $(z_+^*,z_-^*,1)$ attains this sum; if
$\theta_-^*>0$ and $\theta_+^*=0$, the point $(z_+^*,z_-^*,0)$ does. Every
optimal point must attain both base minima, hence equal $z_\pm^*$ in those
coordinates, and must make both added terms vanish, which fixes $y$.
\end{proof}

At the minimizer in (b) the curvature in the $y$ direction is
$2\eta_0\bigl((\theta_+^*)^2+(\theta_-^*)^2\bigr)$. It is positive but can be
extremely small; the amplifier transfers a tiny amplitude into an order-one
coordinate without a usable growth constant.

\paragraph{Square Root Sum.}
An instance consists of positive integers $a_1,\ldots,a_n$ and an integer
$B$; the question is whether $\sum_i\sqrt{a_i}\le B$. If $B<0$ the answer is
no.

\begin{lemma}[Averaging tree]\label{lem:points-tree}
For a Square Root Sum instance with $B\ge0$, compute powers of two $R_i$
with $R_i^2\le a_i<4R_i^2$, and put $c_i=a_i/R_i^2\in[1,4)$,
$M_0=\max_iR_i$, $w_i=R_i/M_0\in(0,1]$, and let $K$ be the least power of
two with $K\ge\max\{2,n\}$. Take a complete binary tree with $K$ leaf
slots. Each real leaf $i\le n$ has a variable $u_i\in[1,2]$ and signal
$w_iu_i$; padded leaves have signal $0$. Each internal node $\nu$ has a
variable $v_\nu\in[0,2]$, its own signal $v_\nu$, and residual
$\varsigma_\nu=v_\nu-\frac12(\text{left signal}+\text{right signal})$. Let
$s$ be the root variable and
\[
 F_0(u,v)=\sum_{i=1}^n\bigl(u_i^3/3-c_iu_i\bigr)+\sum_\nu\varsigma_\nu^2
\]
on the box $Y$ of these variables. Then $\nabla^2F_0\succeq\frac18I$ on $Y$,
and the unique minimizer has root value $s^0=\sum_i\sqrt{a_i}/(KM_0)\in[0,2]$.
The construction takes polynomial time; all coefficients and box widths are
bounded by absolute constants; and the bags
$\{v_\nu\}\cup\{\text{variable children of }\nu\}$, joined along the tree,
form a tree decomposition with bags of size at most three.
\end{lemma}

\begin{proof}
The leaf term has derivative $u_i^2-c_i$ and is minimized on $[1,2]$ at
$u_i=\sqrt{c_i}\in[1,2)$. Setting internal variables to the averages of
their children's signals makes every residual zero and keeps them in
$[0,2]$. This point minimizes every summand, and its root is the average of
the $K$ leaf signals,
$\frac1K\sum_iw_i\sqrt{c_i}=\sum_i\sqrt{a_i}/(KM_0)$.

For the Hessian, order the variables and write the residuals as the rows of
$(I-P)$ applied to a direction $d$: the row of a leaf variable is the
identity row, and the row of an internal variable subtracts half of each
variable child's coefficient ($w_i$ for a leaf, $1$ for an internal node).
Every variable is a child of at most one parent, so distinct rows of $P$
have disjoint supports and $PP^{\mathsf T}$ is diagonal with entries at most
$2\cdot\frac14$. Thus $\norm P\le1/\sqrt2<\frac34$. Since $2u_i\ge2$,
\[
 d^{\mathsf T}\nabla^2F_0d\ge2\sum_id_{u_i}^2+2\sum_\nu(\text{residual
 row of }d)^2=2\norm{(I-P)d}^2\ge2\cdot\tfrac1{16}\norm d^2 .
\]
For the decomposition, a leaf variable occurs only in its parent's bag,
which also covers its unary term; adjacent bags share the child internal
variable; every residual lies in its own bag. The normalization uses
integer square roots of bounded size and binary arithmetic on numbers of
polynomial length, and the tree has fewer than $4n+4$ nodes. Coefficients
are bounded because $c_i<4$, $w_i\le1$, and every residual has coefficients
in $[-1,1]$.
\end{proof}

\begin{lemma}[Integral radical sums]\label{lem:points-trace}
If positive integers $a_1,\ldots,a_n$ satisfy $\sum_i\sqrt{a_i}\in\Z$, then
every $a_i$ is a perfect square.
\end{lemma}

\begin{proof}
Let $\mathcal F=\Q(\sqrt{a_1},\ldots,\sqrt{a_n})$, a number field, and let
$\operatorname{Tr}=\operatorname{Tr}_{\mathcal F/\Q}$. If $a_i$ is not a
square, $\Q(\sqrt{a_i})$ has degree two and
$\operatorname{Tr}_{\Q(\sqrt{a_i})/\Q}(\sqrt{a_i})=0$, so transitivity of the
trace gives $\operatorname{Tr}(\sqrt{a_i})=[\mathcal
F:\Q(\sqrt{a_i})]\cdot0=0$. A rational number $z$ has
$\operatorname{Tr}(z)=[\mathcal F:\Q]z$. Applying the trace to
$\sum_i\sqrt{a_i}=B'\in\Z$ gives $\sum_{a_i\text{ square}}\sqrt{a_i}=B'$, and
subtracting, the sum of the positive numbers $\sqrt{a_i}$ over the
nonsquares is zero. So there are no nonsquares.
\end{proof}

The field $\mathcal F$ appears only in the proof. The reduction tests squareness by
integer square roots.

\begin{proof}[Proof of Proposition~\ref{prop:points-active}]
If $B<0$, answer no. If $B\ge2KM_0$, answer yes, since $s^0\le2$ gives
$\sum_i\sqrt{a_i}=KM_0s^0\le B$. Otherwise $b=B/(KM_0)\in[0,2)$. Apply
Lemma~\ref{lem:points-sign-test} to $\Phi=F_0$ of Lemma~\ref{lem:points-tree}
with $\zeta=b-s$, and let $F=H$. Then $\nabla^2F\succeq\frac1{16}I$, and the
distinguished coordinate $\theta$ satisfies $\theta^*=0$ exactly when
$b\ge s^0$, that is, when $\sum_i\sqrt{a_i}\le B$, including equality;
otherwise $0<\theta^*<1$. The polynomial $F$ is a cubic without constant
term, with coefficients bounded by an absolute constant. Its diagonal
Hessian entries are at most $4+\frac12$ for leaf variables, $2+\frac12$ for
internal variables, and $2$ for $\theta$. Adding the bag $\{s,\theta\}$ at
the root keeps bags of size at most three, so the interaction graph has
treewidth at most two. Point approximation follows from strong convexity:
Lemma~\ref{lem:convex-value} with $\eta=2^{-2q}/32$ returns a feasible point
of gap at most $\eta$, and $F(z)-F_*\ge\frac1{32}\norm{z-z^*}^2$ gives
distance at most $2^{-q}$.
\end{proof}

\begin{proof}[Proof of Proposition~\ref{prop:points-selector}]
Apply Lemma~\ref{lem:points-amplifier}(a) to the cubic $H=F$ of
Proposition~\ref{prop:points-active}, with distinguished coordinate
$\theta$, and a new coordinate $y\in[0,1]$. If $\sum_i\sqrt{a_i}\le B$, then
$\theta^*=0$ and $S=\{z^*\}\times[0,1]$, whose minimum-norm point has $y=0$;
otherwise $S=\{(z^*,1)\}$. The point $(z^*,1)$ is optimal in both cases, and
a rational point within $2^{-q}$ of it is obtained by approximating $z^*$ as
in the proof of Proposition~\ref{prop:points-active} and appending $y=1$.
The new term has scope $\{\theta,y\}$; attaching the bag $\{\theta,y\}$ to
the bag $\{s,\theta\}$ keeps the treewidth at most two. Its coefficients are
bounded.
\end{proof}

\begin{proof}[Proof of Theorem~\ref{thm:points-quartic-lower}(a)]
Handle $B<0$ and $B\ge2KM_0$ as above, and decide the instance by integer
arithmetic if every $a_i$ is a perfect square. Otherwise
Lemma~\ref{lem:points-trace} shows $\sum_i\sqrt{a_i}\ne B$, so $s^0\ne b$.
Build two independent copies of $F_0$ with root variables $s_\pm$, and apply
Lemma~\ref{lem:points-sign-test} with $\zeta_+=b-s_+$ and $\zeta_-=s_--b$.
The resulting $H_\pm$ have Hessians at least $\frac1{16}I$, and
\[
 \theta_+^*>0\iff s^0>b,\qquad\theta_-^*>0\iff s^0<b .
\]
Exactly one amplitude is positive. Lemma~\ref{lem:points-amplifier}(b)
gives a quartic $F$, convex on its box, with a unique minimizer whose
$y$-coordinate is $1$ if $\sum_i\sqrt{a_i}>B$ and $0$ if
$\sum_i\sqrt{a_i}<B$.

For the structural claims, take the two tree decompositions with their
root bags $\{s_\pm,\theta_\pm\}$, add the bags $\{\theta_+,y\}$ and
$\{\theta_-,y\}$ attached to $\{s_+,\theta_+\}$ and $\{s_-,\theta_-\}$, and
join the two new bags. This is a tree decomposition with bags of size at
most three; $y$ occurs in two adjacent bags. Before rescaling, the diagonal
Hessian entries are at most $5$: the amplifier adds at most $2\eta_0$ to the
$\theta_\pm$ entries and contributes $2\eta_0(\theta_+^2+\theta_-^2)$ to the
$y$ entry. The coordinate intervals have width at most two; the affine maps
onto $[0,1]$ preserve degrees and factor scopes and multiply diagonal Hessian
entries by at most four, giving the bound $20$. Each factor (leaf term,
residual square, $\theta^2$, sign test, amplifier) has degree at most four,
at most three variables, and bounded coefficients after the substitution,
and each variable lies in at most three factors. Hence every nonconstant
monomial of the expanded polynomial has a coefficient bounded by an
absolute constant. The constant term is a sum over all factors and need not
be bounded; deleting it does not change the minimizer.
\end{proof}

\paragraph{$\PosSLP$.}
A straight-line program uses the constants $0$ and $1$ and the operations
$+,-,\times$ on earlier results; its output is an integer $A$, and $\PosSLP$
asks whether $A>0$ \citep{AllenderEtAl2009}.

\begin{lemma}[A strongly convex gate objective]\label{lem:points-gates}
From a straight-line program with output $A$ one computes in polynomial time
a rational quartic $\mathcal G$ on the box $[-\frac14,\frac14]^N$ with
$\nabla^2\mathcal G\succeq I$ there, whose unique minimizer $x^*$ has last
coordinate $t^*$ with $t^*>0$ if $A>0$ and $t^*<0$ if $A\le0$. The number
$N$ is linear in the program length, and the exact values of the gates are
never written.
\end{lemma}

\begin{proof}
Represent each gate value $\alpha$ by a pair $(u,v)$ with $v>0$ and
$u/v=\alpha$: the constants $0$ and $1$ by $(0,\frac14)$ and
$(\frac14,\frac14)$, a sum or difference of $(u_1,v_1)$ and $(u_2,v_2)$ by
$\bigl((u_1v_2\pm u_2v_1)/4,\ v_1v_2/4\bigr)$, and a product by
$\bigl(u_1u_2/4,\ v_1v_2/4\bigr)$. These preserve the ratios and the
positivity of $v$, and all entries stay in $[-\frac14,\frac14]$. For the
output pair $(u,v)$ add $t=(2u-v)/4=\frac v4(2A-1)$, which is nonzero and
positive exactly when $A>0$.

List the pair coordinates and $t$ in topological order as $x_1,\ldots,x_N$,
with $x_i=p_i(x_1,\ldots,x_{i-1})$, where each $p_i$ is a constant, an affine
function, or a quadratic with at most two monomials. On the whole box,
repeated operands included,
\[
 |p_i|\le\tfrac14,\qquad\norm{\nabla p_i}_1\le1,\qquad
 \norm{\nabla^2p_i}\le1 .
\]
Let $r_i=x_i-p_i$, $a_i=64^{N-i}$, and $\mathcal G=\sum_ia_ir_i^2$. The
triangular system $r=0$ has the unique solution $x^*$ given by the exact gate
values, so $\mathcal G\ge0=\mathcal G(x^*)$.

Let $D_a=\diag(8^{N-i})$, let $L_p$ be the strictly lower
triangular matrix with entries $\partial_jp_i$, and
$\mathcal E=D_aL_pD_a^{-1}$, with entries $8^{j-i}\partial_jp_i$ for $j<i$.
Row sums of $|\mathcal E|$ are at most $\frac18$ and column sums at most
$\sum_{k\ge1}8^{-k}=\frac17$, so
$\norm{\mathcal E}\le(1/56)^{1/2}<\frac14$. For a direction $h$,
\[
 2\norm{D_a(I-L_p)h}^2=2\norm{(I-\mathcal E)D_ah}^2\ge\tfrac98\norm{D_ah}^2 .
\]
The remaining Hessian term is $-2\sum_ia_ir_ih^{\mathsf T}\nabla^2p_ih$.
Since $|r_i|\le\frac12$ and $p_i$ depends only on earlier coordinates, it is
at least
\[
 -\sum_ia_i\norm{h_{<i}}^2=-\sum_jh_j^2\sum_{i>j}a_i\ge-\tfrac1{63}\norm{D_ah}^2 .
\]
Thus $\nabla^2\mathcal G\succeq(\frac98-\frac1{63})D_a^2\succeq I$. The
polynomial has $O(N)$ monomials and weights of $O(N)$ bits.
\end{proof}

\begin{proof}[Proof of Theorem~\ref{thm:points-quartic-lower}(b)]
Take two copies of $\mathcal G$ with outputs $t_\pm$ and apply
Lemma~\ref{lem:points-sign-test} with $\kappa=1$, $\zeta_+=-t_+$, and
$\zeta_-=t_-$; here $|\zeta_\pm|\le\frac14$. Since $t^*\ne0$, exactly one of
$\theta_+^*>0$ (when $A>0$) and $\theta_-^*>0$ (when $A\le0$) holds.
Lemma~\ref{lem:points-amplifier}(b) gives a quartic, convex on its box, with
a unique minimizer whose $y$-coordinate is $1$ if $A>0$ and $0$ if $A\le0$.

Multiply the whole objective by $64^{-N}$. This preserves the minimizer and
convexity. The gate weights become $64^{-i}$, and every squared residual has
coefficients bounded by an absolute constant even when a monomial occurs in
many residuals, because $\sum_i64^{-i}<\frac1{63}$; the remaining terms are
bounded as well. The affine map from $[-\frac14,\frac14]$ to $[0,1]$ keeps
these bounds, including that of the constant term. The encoding length is
polynomial; the weights have $O(N)$ bits, so bounded magnitude does not mean
bounded bit length. No treewidth bound is claimed.
\end{proof}

\begin{proof}[Proof of Theorem~\ref{thm:points-quartic-lower}(c)]
Given a source instance, run the reduction. If it answers directly, stop.
Otherwise query the point algorithm for a point $\hat z$ with
$\norm{\hat z-z^*}\le\frac14$ and answer according to whether the designated
coordinate of $\hat z$ exceeds $\frac12$; its true value is $0$ or $1$.
Feasibility of $\hat z$ is not used. The running time is polynomial if the
point algorithm is. If the point algorithm is always correct and runs in
expected polynomial time, so does the decision procedure.
\end{proof}

\paragraph{Regularization.}

\begin{proof}[Proof of Proposition~\ref{prop:points-regularization}]
Write $x_0=\frac12$ (a constant) and
$G_n(x)=8\sum_{i=1}^nx_i^2-\sum_{i=1}^nx_{i-1}^2x_i$, so that the $i=1$ term
is $-x_1/4$ and $F_n=G_n+x_n^2(y-1)^2$.

(a) On $[0,1]^n$ the Hessian of $G_n$ is tridiagonal with diagonal entries
$16-2x_{i+1}\ge14$ ($i<n$) and $16$ ($i=n$), and off-diagonal entries
$-2x_i$ in positions $(i,i+1)$. Each absolute off-diagonal row sum is at
most $4$, so $\nabla^2G_n\succeq10I$. For a direction $(\pi,\varsigma)$ in
$(x,y)$, \eqref{eq:points-amplifier-identity} with $\gamma=1$ gives a Hessian
form of at least $(10-6)\norm\pi^2\ge0$, so $F_n$ is convex on the box.

Let $x^*$ minimize $G_n$ on $[0,1]^n$. At a lower bound $x_i=0$ the partial
derivative is $-x_{i-1}^2$, negative once $x_{i-1}>0$; since $x_0=\frac12$,
induction excludes all lower bounds. At an upper bound the derivative is at
least $16-1-2>0$. Hence $x^*$ is interior and
\[
 x_i^*=\frac{(x_{i-1}^*)^2}{16-2x_{i+1}^*}\ (i<n),\qquad
 x_n^*=\frac{(x_{n-1}^*)^2}{16}.
\]
So $0<x_i^*\le(x_{i-1}^*)^2/14$, and with $\omega_i=x_i^*/14$ we get
$\omega_i\le\omega_{i-1}^2$, $\omega_0=\frac1{28}$, hence
$x_n^*\le14\cdot28^{-2^n}\le2^{-2^{n+1}}$ for $n\ge1$. Since
$F_n\ge G_n\ge G_n(x^*)$ with equality at $(x^*,1)$ and $x_n^*>0$, the
point $(x^*,1)$ is the unique minimizer.

(b) The point $(x^*,0)$ has gap $(x_n^*)^2\le2^{-2^{n+2}}$ and distance one.
An error bound $\dist\le\Gamma\,\mathrm{gap}^\vartheta$ at this point gives
$1\le\Gamma(x_n^*)^{2\vartheta}$, so $\log_2\Gamma\ge\vartheta2^{n+2}$.

(c) For $\lambda>0$, the same boundary arguments show that the regularized
minimizer has interior $x$-coordinates. Minimizing over $y$ for fixed $x$
gives $y_\lambda=x_{n,\lambda}^2/(x_{n,\lambda}^2+\lambda)$, and stationarity
in $x$ gives
\[
 x_{i,\lambda}=\frac{x_{i-1,\lambda}^2}{16+2\lambda-2x_{i+1,\lambda}}\
 (i<n),\qquad
 x_{n,\lambda}=\frac{x_{n-1,\lambda}^2}{16+2\lambda+2(1-y_\lambda)^2}.
\]
All denominators are at least $14$, so the bound of (a) applies:
$x_{n,\lambda}\le2^{-2^{n+1}}$. If $y_\lambda\ge\frac12$ then
$\lambda\le x_{n,\lambda}^2\le2^{-2^{n+2}}$. A positive rational
$\lambda=p_\lambda/q_\lambda$ with this property has
$q_\lambda\ge2^{2^{n+2}}$.

(d) Optimality of $z_\lambda$ and $\norm z^2\le n+1$ on the box give
$F_n(z_\lambda)-F_{n,*}\le(n+1)\lambda$. For the point problem,
$F_n(x,1)=G_n(x)$, so $(x^*,1)$ is approximated by applying
Lemma~\ref{lem:convex-value} to the $10$-strongly convex $G_n$ and appending
$y=1$.
\end{proof}

\paragraph{Algebraic output on a conditioned path.}

\begin{proof}[Proof of Proposition~\ref{prop:points-algebraic-output}]
(a) The chain $a_n=1$, $a_i=\sqrt{a_{i+1}+2}$ stays in $[1,2)$ and makes
every square in $\Psi_n$ vanish; conversely $\Psi_n=0$ forces $x_n=1$ and then the
chain. For $x\in[1,2]^n$ and $e=x-a$ the residuals are
\[
 r_i=x_i^2-x_{i+1}-2=(x_i+a_i)e_i-e_{i+1}\ (i<n),\qquad r_n=e_n .
\]
Thus $r=D_r(I-P_r)e$ with $D_r=\diag(x_1+a_1,\ldots,
x_{n-1}+a_{n-1},1)$ and $P_r=D_r^{-1}S_n$, where $S_n$ is the upper shift.
Since $x_i+a_i\ge2$, $P_r$ is a nilpotent weighted shift with
$\norm{P_r}\le\frac12$, so $\norm{(I-P_r)^{-1}}\le2$; also
$\norm{D_r^{-1}}\le1$. Hence $\norm e\le2\norm r$ and
$\Psi_n(x)=\norm r^2\ge\frac14\norm{x-a}^2$. The outgoing factor contributes
$12x_i^2-4x_{i+1}-8\le36$ to $\partial_{ii}\Psi_n$ and the incoming factor
contributes $2$, so $\partial_{ii}\Psi_n\le38$.

(b) Let $\phi(t)=t^2-2$ and $P_k=\phi^{\circ k}-1$, $k=n-1$. Since
$\phi(a_i)=a_{i+1}$, $P_k(a_1)=0$. Modulo $2$, $\phi(t)\equiv t^2$, so
$P_k(t+1)\equiv(t+1)^{2^k}-1\equiv t^{2^k}$. Moreover $\phi(\pm1)=-1$ gives
$P_k(1)=-2$. Hence $P_k(t+1)$ is monic, all its other coefficients are even,
and its constant term is $-2$; Eisenstein's criterion at $2$ shows that it
is irreducible, and so is $P_k$. Therefore $P_k$ is the minimal polynomial of
$a_1$ and $[\Q(a_1):\Q]=2^{n-1}$.

For the support, we show by induction that $\phi^{\circ k}$ ($k\ge1$) has a
nonzero coefficient at every even power $t^{2j}$, $0\le j\le2^{k-1}$, with
sign $(-1)^{2^{k-1}-j}$. This holds for $t^2-2$. If $\chi=\phi^{\circ k}$
has this property, every product contributing to the coefficient of
$t^{2J}$ in $\chi^2$ has sign $(-1)^{2^k-J}$, so these coefficients are
nonzero with that sign. The constant coefficient of $\chi$ is
$\phi^{\circ k}(0)\in\{-2,2\}$, so the constant coefficient of $\chi^2-2$ is
$2$, which has the required sign. Finally the constant coefficient of $P_k$
is $-3$ for $k=1$ and $1$ for $k\ge2$. Hence $P_k$ has exactly
$2^{k-1}+1=2^{n-2}+1$ nonzero coefficients.

(c) Let $\alpha_n=1$ and, for $i=n-1,\ldots,1$, let
$\alpha_i=\lfloor2^8\sqrt{\alpha_{i+1}+2}\rfloor/2^8$, computed by an integer
square root. With $e_i=|\alpha_i-a_i|$, we have $e_n=0$ and
\[
 e_i\le2^{-8}+\frac{e_{i+1}}{\sqrt{\alpha_{i+1}+2}+\sqrt{a_{i+1}+2}}
 \le2^{-8}+\tfrac12e_{i+1},
\]
so $e_i\le2^{-7}$ by induction. Put $\ell_i=\max\{1,\alpha_i-2^{-6}\}$ and
$u_i=\min\{2,\alpha_i+2^{-6}\}$. Then $a\in B=\prod_i[\ell_i,u_i]\subseteq
[1,2]^n$, the endpoints have $O(1)$ bits, and $|x_i-a_i|\le3\cdot2^{-7}$
for $x\in B$. Consequently $|r_i|\le4\cdot3\cdot2^{-7}+3\cdot2^{-7}=
\frac{15}{128}$. The residual Jacobian is
$J=\diag(2x_1,\ldots,2x_{n-1},1)-S_n$; as in (a),
$\norm{J^{-1}}\le2$ because $x_i\ge1$. The exact Hessian is
\[
 \nabla^2\Psi_n=2J^{\mathsf T}J+4\diag(r_1,\ldots,r_{n-1},0)
 \succeq\Bigl(\tfrac12-\tfrac{60}{128}\Bigr)I=\tfrac1{32}I
\]
on $B$.
\end{proof}

\paragraph{Rectangular certificates.}

\begin{proof}[Proof of Proposition~\ref{prop:points-rectangle}]
(a) The residuals are $\varrho_0=x_0-\frac14$ and
$\varrho_i=x_i-\frac14x_{i-1}^2$. They vanish exactly at
$a_0=\frac14$, $a_i=\frac14a_{i-1}^2$, which gives
$a_i=2^{-(4\cdot2^i-2)}$, a point strictly inside $X=[0,\frac12]^{n+1}$; so
$a$ is the unique minimizer, with value $0$. With $e=x-a$ the residuals are
$\varrho=(I-P_\varrho)e$, where $P_\varrho$ is a lower shift with weights
$(x_{i-1}+a_{i-1})/4\le\frac14$. Hence $\norm e\le\frac43\norm\varrho$ and
$\Xi_n\ge\frac9{16}\norm{x-a}^2$. Direct differentiation gives
$\partial_{ii}\Xi_n=2+\frac34x_i^2-x_{i+1}\le\frac{35}{16}$ for $i<n$ and
$\partial_{nn}\Xi_n=2$. With the residual Jacobian $J$, which has unit diagonal
and subdiagonal entries $-x_{i-1}/2\in[-\frac14,0]$, the exact Hessian is
$2J^{\mathsf T}J-\diag(\varrho_1,\ldots,\varrho_n,0)$. Since
$\norm{I-J}\le\frac14$ and $\varrho_i\le\frac12$ on $X$,
$\nabla^2\Xi_n\succeq(2\cdot\frac9{16}-\frac12)I=\frac58I$.

If a rectangle with rational endpoints contains $a$ and lies in
$(0,\frac12)^{n+1}$, its lower endpoint in coordinate $n$ satisfies
$0<\ell_n\le a_n=2^{-m}$ with $m=4\cdot2^n-2$. Writing $\ell_n=p_n/q_n$ with
integers $p_n\ge1$, $q_n\ge1$ gives $q_n\ge2^m$.

(b) Let $h(x)=x_n-\frac18x_{n-1}^2=\frac12(x_n+\varrho_n)$ and
$\widetilde\Xi_n=\Xi_n+y^2+\frac14yh$ on $X\times[0,\frac12]$. Since $x_n,y\ge0$,
$\frac14yh\ge\frac18y\varrho_n\ge-\frac1{16}(y^2+\varrho_n^2)$, and
$\varrho_n^2\le\Xi_n$, so
\[
 \widetilde\Xi_n\ge\tfrac{15}{16}(\Xi_n+y^2)\ge\tfrac{135}{256}
 \bigl(\norm{x-a}^2+y^2\bigr).
\]
Thus $(a,0)$ is the unique minimizer, with value $0$, and
$\partial_y\widetilde\Xi_n(a,0)=\frac14h(a)=\frac18a_n>0$ because
$a_{n-1}^2=4a_n$. The added term changes the diagonal only by
$-\frac y{16}$ in coordinate $n-1$, and $\partial_{yy}\widetilde\Xi_n=2$. Its
other Hessian entries are $-x_{n-1}/16$ in position $(n-1,y)$ and
$\frac14$ in position $(n,y)$; the absolute row sums of this perturbation are
at most $\frac9{32}$, so $\nabla^2\widetilde\Xi_n\succeq(\frac58-\frac9{32})I=
\frac{11}{32}I$. Treewidth two follows from the path bags
$\{x_{i-1},x_i\}$ and the extra bag $\{x_{n-1},x_n,y\}$.

For a rectangle $\prod_i[\ell_i,u_i]\subseteq X$ containing $a$,
$h$ is increasing in $x_n$ and decreasing in $x_{n-1}\ge0$, so the exact
minimum of $\partial_y\widetilde\Xi_n(x,0)=\frac14h(x)$ over the rectangle is
$\frac14(\ell_n-u_{n-1}^2/8)$, attained at a corner. If it is nonnegative,
then $\ell_n\ge u_{n-1}^2/8\ge a_{n-1}^2/8=a_n/2>0$, while containment gives
$\ell_n\le a_n=2^{-m}$. The denominator bound follows as in (a).

(c) The recurrence $a_i=\frac14a_{i-1}^2$ is an arithmetic circuit of size
$O(n)$. The map $T(x)=(\frac14,\frac14x_0^2,\ldots,\frac14x_{n-1}^2)$
sends the closed box $X$ into itself and has Lipschitz constant at most
$\frac14$ there, so the contraction principle on $X$ certifies the unique
fixed point $a$, and $\Xi_n=\norm{x-T(x)}^2\ge0$ certifies optimality. In the
second family, $\varrho_n(a)=0$ gives $h(a)=a_n/2>0$ symbolically, and the
displayed lower bound certifies global optimality. Neither certificate
requires a rectangle in the open box or a sign on a whole rectangle.
\end{proof}
